% =====================================================================
%  Capítulo 15 · Bioinformática estructural
% =====================================================================
\chapter{Bioinformática estructural}\label{cap:15}

\pgfplotsset{
  capquince pdb/.style={curso, width=0.95\linewidth, height=6.2cm,
    xmin=1976, xmax=2025, ymin=0, ymax=19000, stack plots=y, area style,
    xlabel={año de liberación}, ylabel={entradas nuevas por año},
    xtick={1980,1990,2000,2010,2020}, /pgf/number format/1000 sep={},
    legend pos=north west, enlarge x limits=false},
  capquince plddt/.style={curso, width=0.86\linewidth, height=5.4cm,
    xmin=1, xmax=393, ymin=0, ymax=112, xlabel={residuo de p53 humana},
    ylabel={pLDDT}, ytick={0,50,70,90,100}, grid=none, clip=false,
    xtick={1,50,100,150,200,250,300,350,393}},
}

\epigrafe{La conformación nativa está determinada por la totalidad de las interacciones interatómicas y, por lo tanto, por la secuencia de aminoácidos, en un ambiente dado.}{\textcite{c15-anfinsen1973}}

\begin{objetivos}
\begin{itemize}
  \item Describir los niveles de organización de una proteína, la geometría del enlace peptídico y los ángulos diedros $\phi$ y $\psi$, e interpretar un diagrama de Ramachandran construido con estructuras reales.
  \item Comparar la cristalografía de rayos X, la resonancia magnética nuclear y la criomicroscopía electrónica; leer en un archivo del PDB la resolución, el factor $B$ y el método, y trabajar con el formato mmCIF.
  \item Derivar el algoritmo de Kabsch mediante la descomposición en valores singulares, calcular el RMSD y el TM-score entre dos estructuras y explicar por qué miden cosas distintas.
  \item Explicar cómo la coevolución, el aprendizaje profundo y los modelos de lenguaje permitieron predecir estructuras con exactitud experimental (AlphaFold2, RoseTTAFold, ESMFold, AlphaFold3) e interpretar sus medidas de confianza pLDDT y PAE.
  \item Formular el acoplamiento molecular (\ingles{docking}) como un problema de búsqueda y puntuación, ejecutar AutoDock Vina, validar un protocolo por \ingles{redocking} y reconocer sus limitaciones.
\end{itemize}
\end{objetivos}

Durante casi todo este libro las proteínas han sido cadenas de letras. Las hemos alineado, agrupado en familias, anotado y contado. Pero una proteína no actúa como una cadena: actúa como un objeto tridimensional, plegado con una precisión de décimas de nanómetro, cuyas cavidades reconocen un sustrato, cuyas superficies se acoplan con otras proteínas y cuyos movimientos transmiten señales. La secuencia es el plano; la estructura es la máquina. La bioinformática estructural se ocupa de ese paso de una dimensión a tres: cómo se describen, se almacenan y se comparan las estructuras, cómo se predicen a partir de la secuencia y cómo se usan para entender la función y diseñar fármacos.

El capítulo avanza en tres lecciones. En la \cref{sec:15-estructura} construiremos el vocabulario geométrico (el enlace peptídico, los ángulos $\phi/\psi$, las hélices y láminas), veremos cómo se determinan experimentalmente las estructuras, qué contiene el Protein Data Bank y cómo se superponen y comparan dos estructuras con el algoritmo de Kabsch. En la \cref{sec:15-prediccion} recorreremos medio siglo de intentos de predecir la estructura a partir de la secuencia, desde el modelado por homología hasta AlphaFold, y aprenderemos a leer las medidas de confianza que acompañan a cada modelo. En la \cref{sec:15-docking} usaremos estructuras para predecir cómo se une una molécula pequeña a una proteína, el punto de partida del diseño racional de fármacos. El formato PDB y las bases de datos estructurales ya aparecieron en el \cref{cap:02}, y los perfiles y alineamientos múltiples del \cref{cap:04} reaparecerán aquí como la materia prima de la predicción.

% ---------------------------------------------------------------------
\section{Estructura de proteínas y el Protein Data Bank}\label{sec:15-estructura}
% ---------------------------------------------------------------------
\enlacerepo{15.1 · Estructura de proteínas y PDB}

\begin{intuicion}[Un collar de cuentas con bisagras rígidas]
Tome un collar hecho de pequeñas placas planas, unidas por sus esquinas con remaches que solo pueden girar. Cada placa es rígida; toda la flexibilidad del collar está en los remaches. Si lo deja caer sobre la mesa adopta una forma cualquiera, pero si las placas llevaran imanes de distintas intensidades, algunas formas serían mucho más estables que otras, y el collar tendería a ``encontrarlas'' por sí solo. Una proteína es ese collar. Las placas son los enlaces peptídicos, planos por razones químicas; los remaches son los dos enlaces de cada carbono alfa, alrededor de los cuales la cadena puede girar; y los imanes son los puentes de hidrógeno, el efecto hidrofóbico y las interacciones electrostáticas. Describir una estructura, en el fondo, es dar la lista de los ángulos de giro de cada remache.
\end{intuicion}

\subsection{Cuatro niveles de organización}

Desde los trabajos de Linderstrøm-Lang en los años cincuenta se describe la estructura de una proteína en cuatro niveles jerárquicos, que siguen siendo el lenguaje común de la disciplina \parencite{c15-branden1999}:

\begin{itemize}
  \item \textbf{Estructura primaria}: la secuencia de aminoácidos, leída del extremo amino (N) al carboxilo (C). Es lo que codifica el gen y lo que hemos estudiado hasta ahora.
  \item \textbf{Estructura secundaria}: patrones locales y repetitivos del esqueleto, estabilizados por puentes de hidrógeno entre grupos N--H y C=O de la cadena principal. Las dos formas dominantes son la hélice $\alpha$ y la lámina $\beta$; el resto se agrupa en giros y lazos.
  \item \textbf{Estructura terciaria}: la disposición tridimensional completa de una cadena, es decir, cómo se empaquetan hélices y láminas entre sí. Sus unidades de plegamiento independientes se llaman \emph{dominios}, y suelen tener entre 50 y 250 residuos.
  \item \textbf{Estructura cuaternaria}: la asociación de varias cadenas en un complejo, como las cuatro subunidades de la hemoglobina.
\end{itemize}

Esta jerarquía no es solo un recurso didáctico. Refleja la física del plegamiento: la estructura secundaria se forma en microsegundos por interacciones locales, mientras que la terciaria requiere que partes distantes de la secuencia se encuentren. Y refleja la evolución: los dominios se recombinan como piezas de un mecano, y dos proteínas pueden compartir un dominio con el mismo plegamiento aunque sus secuencias hayan divergido hasta ser irreconocibles.

\subsection{El enlace peptídico y los ángulos diedros}

Cada aminoácido aporta al esqueleto tres átomos pesados: el nitrógeno amida (N), el carbono alfa (C$_\alpha$), que porta la cadena lateral R, y el carbono carbonílico (C). Los residuos se unen por un enlace amida entre el C de uno y el N del siguiente, el \emph{enlace peptídico}. Por resonancia entre las formas C=O / C--N y C--O$^-$ / C=N$^+$, este enlace tiene un carácter de doble enlace parcial: mide unos \SI{1,33}{\angstrom}, más corto que un enlace C--N simple (\SI{1,47}{\angstrom}), y no gira libremente. En consecuencia, los seis átomos C$_\alpha$, C, O, N, H y el C$_\alpha$ siguiente quedan en un mismo plano, el \emph{plano peptídico} (\cref{fig:15-peptido}).

\begin{figure}[t]
\centering
% Esqueleto peptídico con los ángulos diedros φ, ψ y ω
\begin{tikzpicture}[font=\sffamily\small,
  at/.style={circle,draw=#1!70!black,fill=#1!80,minimum size=6.4mm,inner sep=0pt,text=white,font=\sffamily\bfseries\scriptsize},
  hat/.style={circle,draw=tinta2,fill=white,minimum size=4.6mm,inner sep=0pt,text=tinta,font=\sffamily\scriptsize},
  enl/.style={line width=1.6pt,tinta2},
  rot/.style={-{Stealth[length=2mm]},thick,#1}]
  % planos peptídicos
  \fill[azulpalido,opacity=.75,rounded corners=4pt] (-0.45,-0.35) -- (1.0,1.35+0.9) -- (2.95,1.35+0.9) -- (4.05,0.95) -- (2.75,-1.55) -- (0.9,-1.55) -- cycle;
  \fill[amarillo!25,opacity=.9,rounded corners=4pt] (3.25,0.95) -- (4.35,-1.8) -- (6.3,-1.8) -- (7.65,-0.3) -- (6.45,2.3) -- (4.5,2.3) -- cycle;
  \node[etiqueta,text=azulprofundo,anchor=west] at (-0.4,2.45) {plano peptídico $i-1$};
  \node[etiqueta,text=amarillo!40!black,anchor=east] at (7.7,2.55) {plano peptídico $i$};
  % enlaces
  \draw[enl] (0,0) -- (1.2,0.7) -- (2.4,0) -- (3.6,0.7) -- (4.8,0) -- (6.0,0.7) -- (7.2,0);
  \draw[enl] (1.2,0.7) -- (1.2,1.85);
  \draw[enl] (1.32,0.75) -- (1.32,1.8);
  \draw[enl] (4.8,0) -- (4.8,-1.15);
  \draw[enl] (4.92,-0.05) -- (4.92,-1.1);
  \draw[enl] (2.4,0) -- (2.4,-1.0);
  \draw[enl] (6.0,0.7) -- (6.0,1.7);
  \draw[enl] (3.6,0.7) -- (3.6,1.9);
  \draw[enl] (3.6,0.7) -- (2.95,1.45);
  \draw[enl,dashed,gris] (1.28,0.58) -- (2.36,-0.06);
  % átomos
  \node[at=gris] at (0,0) {C$_\alpha$};
  \node[at=tinta2] at (1.2,0.7) {C};
  \node[at=rojo] at (1.26,1.95) {O};
  \node[at=azul] at (2.4,0) {N};
  \node[hat] at (2.4,-1.05) {H};
  \node[at=gris] (ca) at (3.6,0.7) {C$_\alpha$};
  \node[at=violeta] at (3.6,2.0) {R};
  \node[hat] at (2.9,1.5) {H};
  \node[at=tinta2] at (4.8,0) {C};
  \node[at=rojo] at (4.86,-1.25) {O};
  \node[at=azul] at (6.0,0.7) {N};
  \node[hat] at (6.0,1.75) {H};
  \node[at=gris] at (7.2,0) {C$_\alpha$};
  % rotaciones
  \draw[rot=naranja] (2.72,0.43) arc[start angle=150,end angle=420,x radius=0.32,y radius=0.16];
  \node[text=naranja!80!black,font=\sffamily\bfseries] at (2.72,0.85) {$\phi$};
  \draw[rot=aqua!80!black] (4.48,0.43) arc[start angle=30,end angle=300,x radius=0.32,y radius=0.16];
  \node[text=aqua!60!black,font=\sffamily\bfseries] at (4.45,0.85) {$\psi$};
  \draw[rot=rojooscuro] (1.54,0.275) arc[start angle=210,end angle=480,x radius=0.3,y radius=0.15];
  \node[text=rojooscuro,font=\sffamily\bfseries] at (1.6,-0.15) {$\omega$};
  % cotas
  \node[etiqueta,anchor=north] at (1.85,-1.55) {C--N: \SI{1,33}{\angstrom}};
  \node[etiqueta,anchor=north] at (1.85,-1.85) {(doble enlace parcial)};
  \node[etiqueta,text=tinta,anchor=west] at (7.6,-0.45) {\ldots};
  \node[etiqueta,text=tinta,anchor=east] at (-0.45,0) {\ldots};
\end{tikzpicture}

\caption[Esqueleto peptídico y ángulos diedros]{Dos planos peptídicos consecutivos (sombreados) articulados en el C$_\alpha$ del residuo $i$. El enlace C--N tiene carácter de doble enlace parcial (línea discontinua) y mantiene coplanares los seis átomos de cada plano; el ángulo $\omega$ en torno a él vale casi siempre \ang{180} (configuración \ingles{trans}). Toda la libertad conformacional del esqueleto se concentra en dos giros por residuo: $\phi$, alrededor del enlace N--C$_\alpha$, y $\psi$, alrededor del enlace C$_\alpha$--C. La cadena lateral R sobresale del C$_\alpha$ y es la que choca estéricamente con el esqueleto cuando $(\phi,\psi)$ toman valores prohibidos.}
\label{fig:15-peptido}
\end{figure}

Si los planos son rígidos, la conformación del esqueleto queda determinada por los ángulos de giro en torno a los dos enlaces restantes de cada C$_\alpha$. Esos giros se miden con ángulos \emph{diedros}.

\begin{definicion}{Ángulo diedro}{diedro}
Dados cuatro átomos consecutivos con posiciones $\mathbf a,\mathbf b,\mathbf c,\mathbf d\in\R^3$, sean $\mathbf b_1=\mathbf b-\mathbf a$, $\mathbf b_2=\mathbf c-\mathbf b$, $\mathbf b_3=\mathbf d-\mathbf c$, y $\mathbf n_1=\mathbf b_1\times\mathbf b_2$, $\mathbf n_2=\mathbf b_2\times\mathbf b_3$ las normales de los dos planos que forman. El ángulo diedro es
\begin{equation}\label{eq:15-diedro}
\theta = \operatorname{atan2}\!\Big( \big(\mathbf n_1\times\mathbf n_2\big)\cdot\frac{\mathbf b_2}{\lVert\mathbf b_2\rVert},\;\; \mathbf n_1\cdot\mathbf n_2 \Big)\in(-\ang{180},\ang{180}].
\end{equation}
Para el residuo $i$: $\phi_i=\theta(\mathrm C_{i-1},\mathrm N_i,\mathrm C_{\alpha,i},\mathrm C_i)$, $\psi_i=\theta(\mathrm N_i,\mathrm C_{\alpha,i},\mathrm C_i,\mathrm N_{i+1})$ y $\omega_i=\theta(\mathrm C_{\alpha,i},\mathrm C_i,\mathrm N_{i+1},\mathrm C_{\alpha,i+1})$.
\end{definicion}

\begin{simbolos}
\simbolo{$\mathbf a,\mathbf b,\mathbf c,\mathbf d$}{Coordenadas cartesianas (en \si{\angstrom}) de cuatro átomos enlazados en secuencia.}
\simbolo{$\mathbf b_1,\mathbf b_2,\mathbf b_3$}{Vectores de enlace; $\mathbf b_2$ es el eje de giro.}
\simbolo{$\mathbf n_1,\mathbf n_2$}{Normales a los planos $(\mathbf a,\mathbf b,\mathbf c)$ y $(\mathbf b,\mathbf c,\mathbf d)$.}
\simbolo{$\operatorname{atan2}(y,x)$}{Arco tangente de dos argumentos: devuelve el ángulo con el signo correcto en los cuatro cuadrantes.}
\simbolo{$\phi_i,\psi_i,\omega_i$}{Diedros del esqueleto del residuo $i$.}
\end{simbolos}

El signo importa: mirando a lo largo del eje $\mathbf b_2$, un diedro positivo corresponde a un giro en sentido horario del enlace lejano respecto al cercano. Con $\omega\approx\ang{180}$ fijo, una proteína de $n$ residuos queda descrita, en primera aproximación, por $2n$ números: sus pares $(\phi_i,\psi_i)$. Es una reducción enorme respecto a las $3\times 8n$ coordenadas cartesianas de sus átomos pesados, y es la razón por la que tantas ideas de este capítulo, desde el diagrama de Ramachandran hasta los métodos de predicción, trabajan en el espacio de los ángulos.

\subsection{El diagrama de Ramachandran}

No todos los pares $(\phi,\psi)$ son posibles. \textcite{c15-ramachandran1963} modelaron los átomos como esferas rígidas con distancias mínimas de contacto y calcularon, para un dipéptido, qué combinaciones de ángulos provocan choques entre los átomos del esqueleto y el C$_\beta$ de la cadena lateral. El resultado, un mapa de regiones permitidas en el plano $(\phi,\psi)$, se conoce desde entonces como \emph{diagrama de Ramachandran}. Lo notable es que se obtuvo \emph{antes} de que existieran suficientes estructuras para comprobarlo: es una predicción de pura estereoquímica.

\begin{figure}[t]
\centering
\includegraphics[width=\linewidth]{cap15/ramachandran.pdf}
\caption[Diagrama de Ramachandran con datos reales]{Diagrama de Ramachandran de \num{19460} residuos de 79 estructuras cristalográficas del PDB con resolución igual o mejor que \SI{1,0}{\angstrom}, coloreados según la estructura secundaria asignada por DSSP. \textbf{Izquierda}: todos los residuos excepto glicina. Se distinguen la región de hélice derecha ($\alpha_R$), la amplia región $\beta$ en el cuadrante superior izquierdo, con la hélice de poliprolina II (PPII) en su borde derecho, y la pequeña isla de hélice izquierda ($\alpha_L$). \textbf{Derecha}: la glicina, sin cadena lateral, ocupa regiones simétricas prohibidas para los demás residuos: el \SI{55}{\percent} de sus residuos tiene $\phi>0$, frente al \SI{3,4}{\percent} del resto. Datos descargados del RCSB y procesados con \archivo{figuras/cap15/generar.py}.}
\label{fig:15-rama}
\end{figure}

La \cref{fig:15-rama} muestra el diagrama construido con estructuras reales de muy alta resolución. Para cada residuo calculamos $(\phi,\psi)$ y su estructura secundaria con DSSP (que veremos enseguida). Tres regiones concentran casi todos los puntos:

\begin{itemize}
  \item \textbf{Hélice derecha} ($\alpha_R$), en torno a $(\phi,\psi)\approx(-\ang{63},-\ang{42})$. En nuestros datos, los residuos asignados a hélice $\alpha$ tienen medias $\bar\phi=-\ang{64,6}$ y $\bar\psi=-\ang{39,8}$.
  \item \textbf{Región $\beta$}, amplia, en el cuadrante $\phi<0$, $\psi>0$. Los residuos en lámina tienen medias $\bar\phi=-\ang{110,5}$, $\bar\psi=\ang{121,0}$. En su borde, cerca de $(-\ang{75},\ang{145})$, se sitúa la hélice de poliprolina II, frecuente en lazos y regiones desordenadas.
  \item \textbf{Hélice izquierda} ($\alpha_L$), cerca de $(\ang{60},\ang{45})$, poco poblada y ocupada sobre todo por glicina, asparagina y aspartato.
\end{itemize}

Dos residuos son especiales. La \emph{glicina}, cuyo ``cadena lateral'' es un solo hidrógeno, no sufre el choque con el C$_\beta$ y puede adoptar conformaciones con $\phi>0$, lo que la convierte en la bisagra preferida de los giros cerrados. La \emph{prolina}, cuya cadena lateral forma un anillo que se cierra sobre el propio nitrógeno, tiene el ángulo $\phi$ casi bloqueado: en nuestros datos, $\bar\phi_{\mathrm{Pro}}=-\ang{67,3}$ con una desviación típica de solo \ang{10,7}. Por eso la prolina rompe hélices y aparece en los extremos de las láminas.

\begin{cuidado}[Un residuo en zona prohibida no es necesariamente un error]
El diagrama de Ramachandran es la primera herramienta de validación de una estructura: un modelo con muchos residuos fuera de las regiones permitidas probablemente está mal construido. Pero una fracción pequeña de residuos atípicos (\ingles{outliers}) aparece incluso en estructuras excelentes, casi siempre por una razón funcional: un residuo tensionado en un sitio activo, o una conformación forzada por la unión de un ligando. La pregunta correcta ante un \ingles{outlier} es si la densidad electrónica lo respalda.
\end{cuidado}

\subsection{Hélices y láminas: geometría de los puentes de hidrógeno}

¿Por qué esas dos regiones y no otras? Porque son las que permiten que los grupos N--H y C=O del esqueleto, que de otro modo quedarían insatisfechos en el interior hidrofóbico de la proteína, formen puentes de hidrógeno entre sí de forma \emph{regular}. En 1951, antes de que se resolviera ninguna estructura de proteína, \textcite{c15-pauling1951} dedujeron la hélice $\alpha$ a partir de tres restricciones: enlace peptídico plano, puentes de hidrógeno lineales entre el C=O del residuo $i$ y el N--H del residuo $i+4$, y equivalencia de todos los residuos. La hélice resultante avanza unos \SI{1,5}{\angstrom} por residuo y tiene cerca de 3,6 residuos por vuelta, lo que da un paso de \SI{5,4}{\angstrom}; todos los N--H apuntan hacia el extremo N y todos los C=O hacia el C, de modo que la hélice entera es un dipolo. En el mismo año, Pauling y Corey propusieron las láminas $\beta$, en las que los puentes se forman \emph{entre} hebras extendidas vecinas, paralelas o antiparalelas.

\begin{historia}[El modelo antes que el dato]
Pauling cuenta que concibió la hélice $\alpha$ en 1948, convaleciente de un resfriado en Oxford, doblando una hoja de papel en la que había dibujado una cadena polipeptídica con las distancias y ángulos correctos. Cuando \textcite{c15-kendrew1958} publicaron el primer modelo tridimensional de una proteína, la mioglobina de cachalote a \SI{6}{\angstrom} de resolución, las varillas de densidad que recorrían la molécula eran, precisamente, hélices $\alpha$. Kendrew describió la estructura como ``más complicada de lo que había predicho cualquier teoría'' y, sin embargo, sus piezas estaban previstas desde hacía una década.
\end{historia}

\subsection{Asignar estructura secundaria: DSSP}

Mirar una estructura y decir ``esto es una hélice'' parece trivial, pero dos expertos pueden discrepar en los extremos. \textcite{c15-kabsch1983} sustituyeron el juicio visual por un algoritmo, DSSP (\ingles{Dictionary of Secondary Structure of Proteins}), que primero identifica los puentes de hidrógeno del esqueleto con un modelo electrostático sencillo y después reconoce patrones en ellos.

\begin{definicion}{Energía de puente de hidrógeno de DSSP}{dssp}
Para un grupo C=O y un grupo N--H del esqueleto, se colocan cargas parciales $\pm q_1$ en C y O, y $\pm q_2$ en N y H. La energía electrostática de interacción es
\begin{equation}\label{eq:15-dssp}
E = f\,q_1 q_2\left(\frac{1}{r_{\mathrm{ON}}}+\frac{1}{r_{\mathrm{CH}}}-\frac{1}{r_{\mathrm{OH}}}-\frac{1}{r_{\mathrm{CN}}}\right),
\end{equation}
y se declara un puente de hidrógeno si $E<-0{,}5$ kcal/mol.
\end{definicion}

\begin{simbolos}
\simbolo{$q_1,\,q_2$}{Cargas parciales: $q_1=0{,}42e$ en el carbonilo y $q_2=0{,}20e$ en la amida.}
\simbolo{$f$}{Factor dimensional, 332 kcal·\si{\angstrom}/(mol·$e^2$), que expresa la energía en kcal/mol.}
\simbolo{$r_{XY}$}{Distancia (en \si{\angstrom}) entre los átomos $X$ e $Y$ de los dos grupos.}
\end{simbolos}

A partir de la matriz de puentes, DSSP define ``giros'' $n$ (un puente entre el C=O de $i$ y el N--H de $i+n$) y ``puentes'' entre hebras. Dos giros de tipo 4 consecutivos forman una hélice $\alpha$ (código \texttt{H}); los de tipo 3 y 5 forman hélices $3_{10}$ (\texttt{G}) y $\pi$ (\texttt{I}); una escalera de puentes entre hebras es una lámina (\texttt{E}), y un puente aislado, \texttt{B}. Los códigos \texttt{T} (giro), \texttt{S} (curvatura) y el espacio en blanco completan el alfabeto. Esta asignación es la ``verdad'' contra la que se evalúan todos los predictores de estructura secundaria y la que usamos para colorear la \cref{fig:15-rama}: en nuestras 79 estructuras, el \SI{30,9}{\percent} de los residuos está en hélice, el \SI{25,8}{\percent} en lámina y el \SI{43,3}{\percent} restante en giros y lazos.

\begin{codigo}[phi\_psi.py]
from Bio.PDB import MMCIFParser, PDBList, PPBuilder
import math

ruta = PDBList().retrieve_pdb_file("1UBQ", file_format="mmCif",
                                   pdir=".")
est = MMCIFParser(QUIET=True).get_structure("ubq", ruta)
for pp in PPBuilder().build_peptides(est[0]):
    for res, (phi, psi) in zip(pp, pp.get_phi_psi_list()):
        if phi is None or psi is None:        # extremos de cadena
            continue
        print(res.get_resname(), res.id[1],
              f"{math.degrees(phi):7.1f} {math.degrees(psi):7.1f}")
\end{codigo}

\subsection{Cómo se determina una estructura}

Ninguna técnica ``fotografía'' una proteína: todas miden una señal indirecta y construyen un modelo atómico compatible con ella. Conocer la señal es indispensable para saber cuánto confiar en cada parte del modelo.

\paragraph{Cristalografía de rayos X} Una proteína purificada se cristaliza, de modo que billones de copias quedan ordenadas en una red periódica. Al iluminar el cristal con rayos X, los electrones dispersan la radiación y las ondas interfieren constructivamente solo en direcciones discretas, que cumplen la ley de Bragg,
\begin{equation}\label{eq:15-bragg}
2\,d\,\sin\theta = \lambda,
\end{equation}
donde $d$ es el espaciado de una familia de planos del cristal, $\theta$ el ángulo de incidencia y $\lambda$ la longitud de onda. Cada reflexión $\mathbf h=(h,k,l)$ tiene un \emph{factor de estructura} $F_{\mathbf h}=|F_{\mathbf h}|e^{i\varphi_{\mathbf h}}$, y la densidad electrónica en el punto $\mathbf x$ de la celda unidad es su transformada de Fourier inversa:
\begin{equation}\label{eq:15-fourier}
\rho(\mathbf x) = \frac{1}{V}\sum_{\mathbf h} |F_{\mathbf h}|\,e^{i\varphi_{\mathbf h}}\,e^{-2\pi i\,\mathbf h\cdot\mathbf x}.
\end{equation}

\begin{simbolos}
\simbolo{$\rho(\mathbf x)$}{Densidad electrónica (electrones/\si{\angstrom^3}) en la posición fraccionaria $\mathbf x$ de la celda.}
\simbolo{$V$}{Volumen de la celda unidad.}
\simbolo{$|F_{\mathbf h}|$}{Amplitud de la reflexión $\mathbf h$; se obtiene de la intensidad medida, $I_{\mathbf h}\propto|F_{\mathbf h}|^2$.}
\simbolo{$\varphi_{\mathbf h}$}{Fase de la reflexión; \emph{no} se mide (el ``problema de la fase'').}
\end{simbolos}

El detector registra intensidades, no fases, y sin fases la \cref{eq:15-fourier} no puede evaluarse. Resolver el problema de la fase (con átomos pesados, dispersión anómala o, cada vez más, reemplazo molecular a partir de un modelo predicho) es el paso crítico. La \emph{resolución} es el menor espaciado $d$ para el que se midieron reflexiones útiles: a \SI{3,5}{\angstrom} se distingue el trazado de la cadena y las hélices; a \SI{2}{\angstrom}, la forma de las cadenas laterales; por debajo de \SI{1,2}{\angstrom}, átomos individuales e incluso algunos hidrógenos. La calidad del ajuste se mide con el factor $R$ y, para evitar el sobreajuste, con el $R_{\text{free}}$, calculado sobre un \SI{5}{\percent} de reflexiones que se excluyen del refinamiento, una forma temprana de validación cruzada.

\paragraph{Resonancia magnética nuclear (RMN)} La proteína se estudia en disolución. Los espectros multidimensionales revelan qué núcleos (normalmente $^1$H, marcados con $^{15}$N y $^{13}$C) están a menos de unos \SI{5}{\angstrom} entre sí (efecto Overhauser nuclear) y qué valores toman ciertos diedros. El resultado no es una estructura sino un conjunto de \emph{restricciones} de distancia y ángulo, y el modelo se entrega como un \emph{ensamble} de 10 a 40 conformaciones que las satisfacen. La RMN está limitada a proteínas de tamaño modesto (en la práctica, hasta unos 30--40 kDa), pero observa la dinámica en disolución, algo que el cristal oculta.

\paragraph{Criomicroscopía electrónica (crio-EM)} Millones de copias de un complejo se congelan en una capa fina de hielo vítreo y se fotografían con un haz de electrones. Cada imagen es una proyección bidimensional ruidosa desde una orientación desconocida; el procesamiento las clasifica, estima sus orientaciones y reconstruye un mapa tridimensional. La resolución se estima comparando dos reconstrucciones independientes, con la curva de correlación de capas de Fourier (FSC), y se informa donde esta cae por debajo de 0,143. Con la llegada de los detectores directos de electrones y nuevos algoritmos, la técnica pasó en pocos años de producir ``manchas'' a mapas con detalle atómico, lo que \textcite{c15-kuhlbrandt2014} bautizó como la ``revolución de la resolución''. La crio-EM no necesita cristales y brilla con complejos grandes, flexibles o de membrana.

\begin{table}[t]
\centering\small
\caption[Métodos experimentales de determinación de estructuras]{Comparación de los tres métodos experimentales principales. Las fracciones del PDB corresponden a las entradas liberadas en 2025 (consulta a la API del RCSB con \archivo{figuras/cap15/generar.py}).}
\label{tab:15-metodos}
\begin{tabularx}{\linewidth}{@{}>{\raggedright\arraybackslash}p{2.0cm}>{\raggedright\arraybackslash}X>{\raggedright\arraybackslash}X>{\raggedright\arraybackslash}X@{}}
\toprule
 & \textbf{Rayos X} & \textbf{RMN} & \textbf{Crio-EM}\\
\midrule
Muestra & cristal & disolución & partículas en hielo vítreo\\
Señal & difracción (amplitudes) & acoplamientos y NOE entre núcleos & proyecciones 2D\\
Tamaño idóneo & cualquiera que cristalice & $\lesssim$ 40 kDa & $\gtrsim$ 50 kDa\\
Medida de calidad & resolución, $R_{\text{free}}$, factor $B$ & restricciones por residuo, dispersión del ensamble & resolución FSC, resolución local\\
Dinámica & indirecta (factor $B$) & directa (ensamble, relajación) & heterogeneidad por clasificación\\
Entradas en 2025 & \num{10083} (\SI{57,6}{\percent}) & \num{225} (\SI{1,3}{\percent}) & \num{7198} (\SI{41,1}{\percent})\\
\bottomrule
\end{tabularx}
\end{table}

\subsection{El Protein Data Bank}

Desde 1971, las estructuras de macromoléculas se depositan en un único archivo público, el Protein Data Bank (PDB). \textcite{c15-berman2000} describieron su reorganización bajo el consorcio RCSB, con procedimientos estandarizados de depósito, validación y distribución, y hoy lo gestiona el consorcio internacional wwPDB. La descripción reciente de \textcite{c15-burley2023} muestra hasta qué punto el archivo ha cambiado: el portal RCSB.org ofrece, junto a las estructuras experimentales, cerca de un millón de modelos calculados por inteligencia artificial, claramente separados de los experimentales.

\begin{figure}[t]
\centering
\begin{tikzpicture}
\begin{axis}[capquince pdb]
  \addplot[draw=azul!80!black, fill=azul!55] table[x=anio,y=xray] {figuras/cap15/pdb_anual.dat} \closedcycle;
  \addlegendentry{rayos X}
  \addplot[draw=aqua!80!black, fill=aqua!55] table[x=anio,y=nmr] {figuras/cap15/pdb_anual.dat} \closedcycle;
  \addlegendentry{RMN en disolución}
  \addplot[draw=naranja!80!black, fill=naranja!60] table[x=anio,y=em] {figuras/cap15/pdb_anual.dat} \closedcycle;
  \addlegendentry{crio-EM}
\end{axis}
\end{tikzpicture}
\caption[Crecimiento del PDB por método]{Entradas liberadas cada año en el PDB según el método experimental, 1976--2025 (áreas apiladas). La cristalografía domina durante cuatro décadas; la RMN alcanza su máximo en torno a 2005 y declina lentamente. A partir de 2014 la crio-EM crece de forma explosiva: pasa del \SI{2,3}{\percent} de las entradas en 2015 al \SI{16,9}{\percent} en 2020 y al \SI{41,1}{\percent} en 2025, mientras que el número de estructuras cristalográficas se estanca. En total se liberaron \num{246212} entradas en el periodo. Datos: API de búsqueda del RCSB, consulta hecha con \archivo{figuras/cap15/generar.py}.}
\label{fig:15-pdb}
\end{figure}

La \cref{fig:15-pdb} resume medio siglo de crecimiento. Cada entrada del PDB tiene un identificador de cuatro caracteres (como \texttt{1UBQ}, la ubiquitina) y contiene mucho más que coordenadas: la secuencia de cada entidad polimérica, los ligandos, las condiciones experimentales, la resolución, las estadísticas de refinamiento, las referencias bibliográficas y un informe de validación.

\paragraph{El formato mmCIF} El formato histórico del PDB, de columnas fijas y 80 caracteres por línea, se describió en el \cref{cap:02}. Sus límites (cinco dígitos para el número de átomo, un solo carácter para la cadena) lo hicieron insuficiente para los grandes complejos de la crio-EM, y desde 2014 el formato de archivo oficial es PDBx/mmCIF. Un archivo mmCIF es una colección de tablas con nombre, cada una precedida de \texttt{loop\_} y de la lista de sus columnas:

\begin{consola}[Fragmento de 1UBQ.cif (tabla \_atom\_site, simplificada)]
loop_
_atom_site.group_PDB
_atom_site.id
_atom_site.type_symbol
_atom_site.label_atom_id
_atom_site.label_comp_id
_atom_site.label_asym_id
_atom_site.label_seq_id
_atom_site.Cartn_x
_atom_site.Cartn_y
_atom_site.Cartn_z
_atom_site.occupancy
_atom_site.B_iso_or_equiv
ATOM 1 N N   MET A 1 27.340 24.430 2.614 1.00  9.67
ATOM 2 C CA  MET A 1 26.266 25.413 2.842 1.00 10.38
ATOM 3 C C   MET A 1 26.913 26.639 3.531 1.00  9.62
\end{consola}

Al ser autodescriptivo, el formato admite tablas nuevas sin romper los programas existentes. En la práctica casi nunca se lee a mano: \cmd{Biopython} (\py{MMCIFParser}), \cmd{gemmi} o \cmd{biotite} lo convierten en objetos jerárquicos de modelo, cadena, residuo y átomo.

\subsection{El factor $B$: cuánto se mueve cada átomo}

La penúltima columna del fragmento anterior, \texttt{occupancy}, indica qué fracción de las moléculas del cristal tiene el átomo en esa posición (menor que 1 si hay conformaciones alternativas). La última, \texttt{B\_iso\_or\_equiv}, es el \emph{factor de desplazamiento atómico} o factor $B$, que mide cuánto se ``emborrona'' la densidad de cada átomo por vibración térmica y por desorden estático entre las copias del cristal.

\begin{teorema}{Factor $B$ y desplazamiento cuadrático medio}{bfactor}
En el modelo isótropo, la contribución de un átomo a la difracción se atenúa por el factor $\exp(-B\sin^2\theta/\lambda^2)$, con
\begin{equation}\label{eq:15-bfactor}
B = 8\pi^2\,\langle u_x^2\rangle
\qquad\Longrightarrow\qquad
\sqrt{\langle u^2\rangle} = \sqrt{3\langle u_x^2\rangle}=\sqrt{\frac{3B}{8\pi^2}}.
\end{equation}
\end{teorema}

\begin{simbolos}
\simbolo{$B$}{Factor $B$ del átomo, en \si{\angstrom^2}.}
\simbolo{$\langle u_x^2\rangle$}{Desplazamiento cuadrático medio del átomo a lo largo de una dirección.}
\simbolo{$\sqrt{\langle u^2\rangle}$}{Desplazamiento cuadrático medio total en tres dimensiones (isótropo).}
\end{simbolos}

Así, $B=\SI{10}{\angstrom^2}$ corresponde a un desplazamiento tridimensional típico de \SI{0,62}{\angstrom}; $B=20$, a \SI{0,87}{\angstrom}; $B=80$, a \SI{1,74}{\angstrom}. Un átomo con $B$ alto está mal definido en el modelo, y conviene no apoyar conclusiones finas (una distancia de puente de hidrógeno, la orientación de una cadena lateral) en él.

\begin{ejemplo}{Cristal frente a disolución: la ubiquitina}{ubq}
La ubiquitina humana, de 76 residuos, se ha resuelto por cristalografía a \SI{1,8}{\angstrom} (\texttt{1UBQ}) y por RMN con un ensamble de 10 modelos (\texttt{1D3Z}). Superponiendo los C$_\alpha$ con el algoritmo de Kabsch (que veremos a continuación), cada modelo de RMN difiere del cristal en un RMSD medio de \SI{0,97}{\angstrom} (entre 0,47 y 1,55). Si se excluye la cola C-terminal (residuos 71--76), la diferencia cae a \SI{0,37}{\angstrom}: el núcleo es prácticamente idéntico en el cristal y en disolución.

La \cref{fig:15-ubq} explica por qué. El factor $B$ medio de los C$_\alpha$ del cristal es de \SI{8,9}{\angstrom^2} en los residuos 1--70 y de \SI{32,8}{\angstrom^2} en los residuos 72--76, y la fluctuación (RMSF) de esos mismos residuos entre los modelos de RMN alcanza \SI{6,6}{\angstrom} en el último. Dos experimentos completamente distintos coinciden en que la cola es flexible: la correlación entre el factor $B$ y la RMSF es de 0,71. Y no es un detalle menor: la glicina 76 de esa cola es la que se une covalentemente a las proteínas marcadas para su degradación, una función que exige flexibilidad.
\end{ejemplo}

\begin{figure}[t]
\centering
\begin{tikzpicture}
\begin{groupplot}[group style={group size=1 by 2, vertical sep=0.35cm, xticklabels at=edge bottom},
   curso, width=\linewidth, height=3.7cm, xmin=0.5, xmax=76.5, grid=none,
   xtick={1,10,20,30,40,50,60,70,76}]
\nextgroupplot[ylabel={$B$ de C$_\alpha$ (\si{\angstrom^2})}, ymin=0, ymax=45]
  \fill[amarillo!20] (axis cs:70.5,0) rectangle (axis cs:76.5,45);
  \addplot[ybar, bar width=1.9pt, fill=azul!70, draw=none] table[x=res,y=B] {figuras/cap15/ubq.dat};
  \node[etiqueta,anchor=north east,text=amarillo!30!black] at (axis cs:70,44) {cola C-terminal $\rightarrow$};
  \node[etiqueta,anchor=north west] at (axis cs:2,44) {cristal \texttt{1UBQ} (\SI{1,8}{\angstrom})};
\nextgroupplot[ylabel={RMSF (\si{\angstrom})}, xlabel={residuo}, ymin=0, ymax=7.2]
  \fill[amarillo!20] (axis cs:70.5,0) rectangle (axis cs:76.5,7.2);
  \addplot[naranja, very thick, mark=*, mark size=0.9pt] table[x=res,y=rmsf] {figuras/cap15/ubq.dat};
  \node[etiqueta,anchor=north west] at (axis cs:2,6.9) {ensamble de RMN \texttt{1D3Z} (10 modelos)};
\end{groupplot}
\end{tikzpicture}
\caption[Flexibilidad de la ubiquitina medida de dos formas]{Dos medidas independientes de la flexibilidad de la ubiquitina. \textbf{Arriba}: factor $B$ de los C$_\alpha$ en el cristal \texttt{1UBQ}. \textbf{Abajo}: fluctuación cuadrática media (RMSF) de los C$_\alpha$ entre los 10 modelos del ensamble de RMN \texttt{1D3Z}, superpuestos sobre el núcleo (residuos 1--70). Ambas señalan la cola C-terminal (sombreada) como la región más móvil, y ambas detectan una flexibilidad moderada en el lazo de los residuos 7--10. Lo que en cristalografía aparece como densidad borrosa, en RMN aparece como desacuerdo entre modelos.}
\label{fig:15-ubq}
\end{figure}

\subsection{Superponer dos estructuras: el algoritmo de Kabsch}

Una pregunta aparece una y otra vez en bioinformática estructural: ¿cuánto se parecen dos estructuras? Dos modelos de la misma proteína, un modelo predicho y la estructura experimental, una proteína antes y después de unir un ligando. Si tenemos una correspondencia entre átomos (el C$_\alpha$ $k$ de una con el C$_\alpha$ $k$ de la otra), la medida natural es la raíz de la desviación cuadrática media (RMSD). Pero las coordenadas de dos archivos están en sistemas de referencia arbitrarios; antes de medir hay que \emph{superponer}, es decir, encontrar la rotación y la traslación que mejor llevan una estructura sobre la otra.

\begin{definicion}{RMSD tras superposición óptima}{rmsd}
Sean $\{\mathbf p_k\}_{k=1}^N$ y $\{\mathbf q_k\}_{k=1}^N$ dos conjuntos de puntos en $\R^3$ con correspondencia conocida. El RMSD tras superposición óptima es
\begin{equation}\label{eq:15-rmsd}
\mathrm{RMSD} = \min_{\mathbf R\in SO(3),\ \mathbf t\in\R^3}\ \sqrt{\frac1N\sum_{k=1}^{N}\big\lVert \mathbf R\mathbf p_k+\mathbf t-\mathbf q_k\big\rVert^2}.
\end{equation}
\end{definicion}

\begin{simbolos}
\simbolo{$\mathbf p_k,\ \mathbf q_k$}{Coordenadas del átomo $k$ en la estructura móvil y en la de referencia.}
\simbolo{$N$}{Número de pares de átomos comparados (por ejemplo, los C$_\alpha$ alineados).}
\simbolo{$\mathbf R$}{Matriz de rotación: ortogonal ($\mathbf R^\top\mathbf R=\mathbf I$) y con $\det\mathbf R=+1$.}
\simbolo{$SO(3)$}{Grupo de las rotaciones propias en tres dimensiones (sin reflexiones).}
\simbolo{$\mathbf t$}{Vector de traslación.}
\end{simbolos}

\textcite{c15-kabsch1976} mostró que este problema de optimización, que parece requerir una búsqueda sobre todas las rotaciones posibles, tiene solución cerrada. Presentamos la derivación moderna, basada en la descomposición en valores singulares (SVD).

\paragraph{Paso 1: la traslación óptima une los centroides} Fijada $\mathbf R$, la función $E(\mathbf t)=\sum_k\lVert\mathbf R\mathbf p_k+\mathbf t-\mathbf q_k\rVert^2$ es cuadrática en $\mathbf t$. Igualando su gradiente a cero,
\[
\nabla_{\mathbf t}E = 2\sum_k(\mathbf R\mathbf p_k+\mathbf t-\mathbf q_k)=\mathbf 0
\quad\Longrightarrow\quad
\mathbf t^\ast = \bar{\mathbf q}-\mathbf R\bar{\mathbf p},
\]
donde $\bar{\mathbf p}$ y $\bar{\mathbf q}$ son los centroides. Por lo tanto, basta centrar ambas estructuras en el origen, $\mathbf p_k\leftarrow\mathbf p_k-\bar{\mathbf p}$ y $\mathbf q_k\leftarrow\mathbf q_k-\bar{\mathbf q}$, y buscar solo la rotación.

\paragraph{Paso 2: minimizar el error equivale a maximizar una traza} Con los puntos centrados, y usando que una rotación conserva las normas ($\lVert\mathbf R\mathbf p_k\rVert=\lVert\mathbf p_k\rVert$),
\begin{equation}\label{eq:15-traza}
E(\mathbf R)=\sum_k\lVert\mathbf p_k\rVert^2+\sum_k\lVert\mathbf q_k\rVert^2-2\sum_k\mathbf q_k^\top\mathbf R\,\mathbf p_k
=\text{const}-2\,\operatorname{tr}\!\big(\mathbf R\,\mathbf H\big),
\qquad \mathbf H=\sum_k \mathbf p_k\mathbf q_k^\top.
\end{equation}
El paso intermedio usa la identidad $\mathbf q^\top\mathbf R\mathbf p=\operatorname{tr}(\mathbf R\,\mathbf p\,\mathbf q^\top)$. Minimizar $E$ equivale, pues, a maximizar $\operatorname{tr}(\mathbf R\mathbf H)$, donde $\mathbf H$ es la matriz $3\times3$ de covarianza cruzada entre las dos nubes de puntos.

\paragraph{Paso 3: la SVD resuelve la maximización} Escribamos $\mathbf H=\mathbf U\boldsymbol\Sigma\mathbf V^\top$, con $\mathbf U,\mathbf V$ ortogonales y $\boldsymbol\Sigma=\operatorname{diag}(\sigma_1,\sigma_2,\sigma_3)$, $\sigma_1\ge\sigma_2\ge\sigma_3\ge0$. Por la propiedad cíclica de la traza,
\[
\operatorname{tr}(\mathbf R\mathbf U\boldsymbol\Sigma\mathbf V^\top)=\operatorname{tr}(\underbrace{\mathbf V^\top\mathbf R\mathbf U}_{\mathbf M}\boldsymbol\Sigma)=\sum_{j=1}^3 M_{jj}\,\sigma_j .
\]
Como $\mathbf M$ es ortogonal, sus elementos diagonales cumplen $M_{jj}\le1$, así que la suma es máxima cuando $\mathbf M=\mathbf I$, es decir, $\mathbf R=\mathbf V\mathbf U^\top$. Queda un detalle: si $\det(\mathbf V\mathbf U^\top)=-1$, esa matriz es una \emph{reflexión}, que convertiría una proteína en su imagen especular. La mejor rotación propia se obtiene entonces sacrificando el término del valor singular más pequeño.

\begin{teorema}{Algoritmo de Kabsch}{kabsch}
Con las estructuras centradas y $\mathbf H=\sum_k\mathbf p_k\mathbf q_k^\top=\mathbf U\boldsymbol\Sigma\mathbf V^\top$, la rotación óptima y el error mínimo son
\begin{gather}
\mathbf R^\ast=\mathbf V\,\mathbf D\,\mathbf U^\top,\qquad
\mathbf D=\operatorname{diag}\big(1,\,1,\,d\big),\qquad d=\operatorname{sign}\det(\mathbf V\mathbf U^\top),\label{eq:15-kabsch}\\
N\cdot\mathrm{RMSD}^2=\sum_k\lVert\mathbf p_k\rVert^2+\sum_k\lVert\mathbf q_k\rVert^2-2(\sigma_1+\sigma_2+d\,\sigma_3).\nonumber
\end{gather}
\end{teorema}

\begin{simbolos}
\simbolo{$\mathbf H$}{Matriz de covarianza cruzada $3\times3$ entre las estructuras centradas.}
\simbolo{$\mathbf U,\ \mathbf V$}{Matrices ortogonales de vectores singulares izquierdos y derechos de $\mathbf H$.}
\simbolo{$\sigma_j$}{Valores singulares de $\mathbf H$, ordenados de mayor a menor.}
\simbolo{$d$}{Corrección de quiralidad: $+1$ si $\mathbf V\mathbf U^\top$ ya es una rotación, $-1$ si es una reflexión.}
\end{simbolos}

El coste es el de construir $\mathbf H$, $O(N)$, más una SVD de $3\times3$, que es constante: superponer dos proteínas de mil residuos toma microsegundos. La \cref{fig:15-kabsch} muestra los tres pasos en dos dimensiones, donde las matrices son $2\times2$ y la geometría se ve de un vistazo.

\begin{figure}[t]
\centering
\resizebox{\linewidth}{!}{% Superposición de Kabsch en 2D con puntos de generar.py
\begin{tikzpicture}[font=\sffamily\small,
  pq/.style={circle,fill=azul,draw=white,minimum size=2.4mm,inner sep=0pt},
  pp/.style={circle,fill=naranja,draw=white,minimum size=2.4mm,inner sep=0pt},
  corr/.style={gris,densely dotted,thin}]
  \begin{scope}[shift={(0,0)},scale=0.62]
    \coordinate (q0) at (0.000,0.000);
\coordinate (q1) at (1.500,0.300);
\coordinate (q2) at (2.400,1.600);
\coordinate (q3) at (1.200,2.500);
\coordinate (q4) at (-0.300,1.400);
\coordinate (p0) at (5.603,-1.121);
\coordinate (p1) at (6.150,0.366);
\coordinate (p2) at (5.790,1.846);
\coordinate (p3) at (4.222,1.354);
\coordinate (p4) at (4.275,-0.385);
\coordinate (c0) at (1.355,-0.373);
\coordinate (c1) at (1.902,1.114);
\coordinate (c2) at (1.542,2.594);
\coordinate (c3) at (-0.026,2.102);
\coordinate (c4) at (0.027,0.363);
\coordinate (a0) at (-0.034,-0.071);
\coordinate (a1) at (1.485,0.376);
\coordinate (a2) at (2.458,1.548);
\coordinate (a3) at (1.128,2.513);
\coordinate (a4) at (-0.236,1.433);
\def\capquinceRantes{4.48}\def\capquinceRcentro{1.21}\def\capquinceRfinal{0.08}\def\capquinceAng{53}

    \draw[azul!60] (q0)--(q1)--(q2)--(q3)--(q4)--cycle;
    \draw[naranja!70] (p0)--(p1)--(p2)--(p3)--(p4)--cycle;
    \foreach \i in {0,...,4}{\draw[corr] (q\i) -- (p\i);}
    \foreach \i in {0,...,4}{\node[pq] at (q\i) {}; \node[pp] at (p\i) {};}
  \end{scope}
  \begin{scope}[shift={(6.0,0)},scale=0.85]
    \coordinate (q0) at (0.000,0.000);
\coordinate (q1) at (1.500,0.300);
\coordinate (q2) at (2.400,1.600);
\coordinate (q3) at (1.200,2.500);
\coordinate (q4) at (-0.300,1.400);
\coordinate (p0) at (5.603,-1.121);
\coordinate (p1) at (6.150,0.366);
\coordinate (p2) at (5.790,1.846);
\coordinate (p3) at (4.222,1.354);
\coordinate (p4) at (4.275,-0.385);
\coordinate (c0) at (1.355,-0.373);
\coordinate (c1) at (1.902,1.114);
\coordinate (c2) at (1.542,2.594);
\coordinate (c3) at (-0.026,2.102);
\coordinate (c4) at (0.027,0.363);
\coordinate (a0) at (-0.034,-0.071);
\coordinate (a1) at (1.485,0.376);
\coordinate (a2) at (2.458,1.548);
\coordinate (a3) at (1.128,2.513);
\coordinate (a4) at (-0.236,1.433);
\def\capquinceRantes{4.48}\def\capquinceRcentro{1.21}\def\capquinceRfinal{0.08}\def\capquinceAng{53}

    \draw[azul!60] (q0)--(q1)--(q2)--(q3)--(q4)--cycle;
    \draw[naranja!70] (c0)--(c1)--(c2)--(c3)--(c4)--cycle;
    \foreach \i in {0,...,4}{\draw[corr] (q\i) -- (c\i);}
    \foreach \i in {0,...,4}{\node[pq] at (q\i) {}; \node[pp] at (c\i) {};}
  \end{scope}
  \begin{scope}[shift={(10.4,0)},scale=0.85]
    \coordinate (q0) at (0.000,0.000);
\coordinate (q1) at (1.500,0.300);
\coordinate (q2) at (2.400,1.600);
\coordinate (q3) at (1.200,2.500);
\coordinate (q4) at (-0.300,1.400);
\coordinate (p0) at (5.603,-1.121);
\coordinate (p1) at (6.150,0.366);
\coordinate (p2) at (5.790,1.846);
\coordinate (p3) at (4.222,1.354);
\coordinate (p4) at (4.275,-0.385);
\coordinate (c0) at (1.355,-0.373);
\coordinate (c1) at (1.902,1.114);
\coordinate (c2) at (1.542,2.594);
\coordinate (c3) at (-0.026,2.102);
\coordinate (c4) at (0.027,0.363);
\coordinate (a0) at (-0.034,-0.071);
\coordinate (a1) at (1.485,0.376);
\coordinate (a2) at (2.458,1.548);
\coordinate (a3) at (1.128,2.513);
\coordinate (a4) at (-0.236,1.433);
\def\capquinceRantes{4.48}\def\capquinceRcentro{1.21}\def\capquinceRfinal{0.08}\def\capquinceAng{53}

    \draw[azul!60] (q0)--(q1)--(q2)--(q3)--(q4)--cycle;
    \draw[naranja!70] (a0)--(a1)--(a2)--(a3)--(a4)--cycle;
    \foreach \i in {0,...,4}{\draw[corr] (q\i) -- (a\i);}
    \foreach \i in {0,...,4}{\node[pq] at (q\i) {}; \node[pp] at (a\i) {};}
  \end{scope}
  \node[font=\sffamily\bfseries\small,text=tinta] at (2.1,3.3) {(a) posición inicial};
  \node[font=\sffamily\bfseries\small,text=tinta] at (6.9,3.3) {(b) centroides unidos};
  \node[font=\sffamily\bfseries\small,text=tinta] at (11.6,3.3) {(c) rotación óptima};
  \node[etiqueta] at (2.1,-1.25) {RMSD $=\SI{4,48}{\angstrom}$};
  \node[etiqueta] at (6.9,-1.25) {RMSD $=\SI{1,21}{\angstrom}$};
  \node[etiqueta] at (11.4,-1.25) {RMSD $=\SI{0,08}{\angstrom}$ (giro de \ang{-53,4})};
  \draw[flecha=tinta2] (4.35,1.0) -- node[above,etiqueta]{$-\bar{\mathbf p}+\bar{\mathbf q}$} (5.3,1.0);
  \draw[flecha=tinta2] (8.6,1.0) -- node[above,etiqueta]{$\mathbf R=\mathbf V\mathbf D\mathbf U^{\top}$} (9.6,1.0);
  \node[pq,label={[etiqueta]right:referencia $\mathbf q_k$}] at (3.2,-1.9) {};
  \node[pp,label={[etiqueta]right:estructura móvil $\mathbf p_k$}] at (6.2,-1.9) {};
\end{tikzpicture}
}
\caption[Superposición de Kabsch]{El algoritmo de Kabsch en dos dimensiones, con cinco pares de puntos. (a) Las estructuras están en posiciones y orientaciones arbitrarias; las líneas punteadas unen los puntos correspondientes. (b) Al trasladar el centroide de la estructura móvil (naranja) sobre el de la referencia (azul), el RMSD cae de 4,48 a \SI{1,21}{\angstrom}. (c) La rotación $\mathbf R=\mathbf V\mathbf D\mathbf U^\top$ obtenida de la SVD de $\mathbf H$ elimina el resto, salvo el ruido añadido a propósito a cada punto: RMSD $=\SI{0,08}{\angstrom}$. Coordenadas generadas con \archivo{figuras/cap15/generar.py}.}
\label{fig:15-kabsch}
\end{figure}

\begin{ejemplo}{Kabsch a mano, en el plano}{kabsch}
En la \cref{fig:15-kabsch}, la estructura móvil se construyó girando la de referencia \ang{52} alrededor de su centroide, trasladándola y añadiendo ruido de unas centésimas. Tras centrar, las sumas de cuadrados son $\sum\lVert\mathbf p_k\rVert^2=9{,}154$ y $\sum\lVert\mathbf q_k\rVert^2=9{,}064$, y la matriz de covarianza cruzada vale
\[
\mathbf H=\begin{pmatrix}1{,}906 & -2{,}558\\ 4{,}741 & 3{,}519\end{pmatrix},
\qquad \sigma_1=5{,}905,\quad \sigma_2=3{,}190,\quad d=+1.
\]
La \cref{eq:15-kabsch} (en su versión bidimensional, con $\mathbf D=\operatorname{diag}(1,d)$) da $N\cdot\mathrm{RMSD}^2=9{,}154+9{,}064-2(5{,}905+3{,}190)=0{,}029$, de modo que $\mathrm{RMSD}=\sqrt{0{,}029/5}=\SI{0,076}{\angstrom}$, sin haber calculado todavía la rotación. La matriz $\mathbf R^\ast=\mathbf V\mathbf U^\top$ resulta ser $\left(\begin{smallmatrix}0{,}597 & 0{,}803\\ -0{,}803 & 0{,}597\end{smallmatrix}\right)$, un giro de $-\ang{53,4}$: el algoritmo recupera el giro original de \ang{52} con un error de algo más de un grado, causado por el ruido.
\end{ejemplo}

\begin{codigo}[kabsch.py]
import numpy as np

def kabsch_rmsd(P, Q):
    """P, Q: arrays (N, 3) con correspondencia fila a fila."""
    P = P - P.mean(axis=0)                 # paso 1: centrar
    Q = Q - Q.mean(axis=0)
    H = P.T @ Q                            # paso 2: covarianza 3x3
    U, S, Vt = np.linalg.svd(H)            # paso 3: SVD
    d = np.sign(np.linalg.det(Vt.T @ U.T))  # evitar reflexiones
    D = np.diag([1.0, 1.0, d])
    R = Vt.T @ D @ U.T
    diff = P @ R.T - Q
    return np.sqrt((diff ** 2).sum(axis=1).mean()), R
\end{codigo}

\subsection{Más allá del RMSD: TM-score}

El RMSD tiene dos defectos serios como medida de similitud global. Primero, depende de la longitud: dos proteínas grandes no relacionadas tienen un RMSD mayor que dos pequeñas, así que un mismo valor no significa lo mismo para 50 que para 500 residuos. Segundo, está dominado por los peores pares: como promedia \emph{cuadrados}, una cola desordenada o un dominio que gira como bisagra puede inflar el RMSD de una estructura que, en el resto, es idéntica. La ubiquitina lo ilustró: \SI{0,97}{\angstrom} con la cola, \SI{0,37}{\angstrom} sin ella.

\textcite{c15-zhang2004} propusieron una medida que resuelve ambos problemas.

\begin{definicion}{TM-score}{tmscore}
Para una estructura modelo y una de referencia de longitud $L$, con $N_{\mathrm{al}}$ pares de residuos alineados,
\begin{equation}\label{eq:15-tmscore}
\mathrm{TM} = \max_{\mathbf R,\mathbf t}\ \frac1L\sum_{k=1}^{N_{\mathrm{al}}}\frac{1}{1+\left(d_k/d_0(L)\right)^2},
\qquad d_0(L)=1{,}24\sqrt[3]{L-15}-1{,}8,
\end{equation}
donde $d_k$ es la distancia entre el par $k$ tras la superposición $(\mathbf R,\mathbf t)$.
\end{definicion}

\begin{simbolos}
\simbolo{$L$}{Longitud de la estructura de referencia (normaliza el puntaje).}
\simbolo{$N_{\mathrm{al}}$}{Número de pares de residuos alineados.}
\simbolo{$d_k$}{Distancia (\si{\angstrom}) entre los C$_\alpha$ del par $k$ tras la superposición.}
\simbolo{$d_0(L)$}{Escala de distancia que crece con la longitud; para $L=141$ vale \SI{4,42}{\angstrom}.}
\end{simbolos}

Cada par contribuye entre 0 y 1: cerca de 1 si $d_k\ll d_0$ y casi nada si $d_k\gg d_0$, de modo que los residuos muy desviados dejan de pesar en lugar de dominar. La escala $d_0(L)$ se eligió para que el puntaje medio de estructuras no relacionadas apenas dependa de la longitud. El resultado va de 0 a 1; en la práctica, un TM-score mayor que 0,5 indica, casi siempre, el mismo plegamiento. Observe también el ``máx'': la superposición óptima para el TM-score \emph{no} es la de Kabsch (que minimiza el RMSD), sino la que maximiza el número de residuos bien superpuestos, y se busca de forma heurística a partir de fragmentos semilla.

\begin{ejemplo}{Mioglobina frente a hemoglobina}{globinas}
La mioglobina de cachalote (\texttt{1A6M}, 151 residuos) y la cadena $\alpha$ de la hemoglobina humana (\texttt{2DN2}, 141 residuos) son un ejemplo clásico de homología estructural. Un alineamiento global con BLOSUM62 empareja 138 residuos con solo un \SI{28,3}{\percent} de identidad, cerca de la ``zona crepuscular'' del \cref{cap:03}. Sin embargo, la superposición de Kabsch de esos 138 C$_\alpha$ da un RMSD de \SI{1,93}{\angstrom}; si se descartan iterativamente los pares a más de \SI{3}{\angstrom}, quedan 129 con un RMSD de \SI{1,38}{\angstrom}. El TM-score, normalizado por la cadena $\alpha$, es 0,87, inequívocamente el mismo plegamiento. Como control, si se baraja al azar la correspondencia entre residuos, el TM-score cae a 0,12. \emph{La estructura se conserva mucho más que la secuencia}, una regla que justifica el modelado por homología de la próxima lección.
\end{ejemplo}

\begin{ideaclave}
El RMSD responde ``¿cuán lejos están los átomos, en promedio?'' y lo dominan los peores; el TM-score responde ``¿qué fracción de la estructura coincide?'' y es comparable entre proteínas de distinta longitud.
\end{ideaclave}

\subsection{Clasificar el universo de plegamientos: SCOP y CATH}

Si la estructura se conserva más que la secuencia, entonces clasificar las proteínas por su estructura revela parentescos que la secuencia ya no muestra. \textcite{c15-murzin1995} crearon SCOP (\ingles{Structural Classification of Proteins}), una clasificación jerárquica hecha en gran parte a mano por expertos, con cuatro niveles: \emph{clase} (todo $\alpha$, todo $\beta$, $\alpha/\beta$ con hebras y hélices alternadas, $\alpha+\beta$ con regiones separadas), \emph{plegamiento} (misma disposición y topología de los elementos secundarios), \emph{superfamilia} (probable origen evolutivo común, inferido por estructura y función aunque la identidad de secuencia sea baja) y \emph{familia} (homología clara a nivel de secuencia). Casi al mismo tiempo, \textcite{c15-orengo1997} presentaron CATH, más automatizada, cuyo nombre enumera sus niveles: \emph{Clase}, \emph{Arquitectura} (forma general, sin considerar la conectividad), \emph{Topología} (conectividad, equivalente al plegamiento) y superfamilia \emph{Homóloga}. Ambas clasifican \emph{dominios}, no cadenas completas, y ambas mostraron que el número de plegamientos distintos en la naturaleza es sorprendentemente pequeño: unos pocos miles, reutilizados una y otra vez.

\begin{profundizacion}[Alineamiento estructural sin correspondencia conocida]
Todo lo anterior supone que sabemos qué residuo corresponde a cuál. Cuando comparamos dos proteínas distintas, esa correspondencia es precisamente lo que queremos averiguar, y el problema se vuelve mucho más difícil: hay que optimizar simultáneamente el alineamiento y la superposición. Programas como TM-align alternan ambos pasos (alinear con programación dinámica sobre una matriz de distancias tras superponer; superponer con los pares alineados), y herramientas recientes como Foldseek traducen la geometría local a un ``alfabeto estructural'' de 20 letras para usar la maquinaria rápida de búsqueda de secuencias del \cref{cap:03} sobre millones de estructuras.
\end{profundizacion}

% ---------------------------------------------------------------------
\section{Predicción de estructura: de la homología a AlphaFold y ESMFold}\label{sec:15-prediccion}
% ---------------------------------------------------------------------
\enlacerepo{15.2 · Predicción de estructura: AlphaFold y ESMFold}

\begin{intuicion}[Reconstruir un edificio a partir de sus reformas]
Suponga que quiere saber la forma de un edificio antiguo que nunca ha visto, pero dispone de los registros de reformas de cientos de edificios hermanos construidos con el mismo plano. Observa algo curioso: cada vez que alguien cambió una viga en la tercera planta, casi siempre cambió también un pilar de la séptima, en otro ala. La explicación más sencilla es que la viga y el pilar se \emph{tocan}: modificar uno obliga a ajustar el otro. Con suficientes registros podría reconstruir qué partes están en contacto y, de ahí, la forma del edificio, sin haberlo medido nunca. Así funciona la predicción de estructura moderna: las familias de proteínas son los edificios hermanos, las mutaciones son las reformas, y las mutaciones correlacionadas delatan qué residuos están juntos en el espacio.
\end{intuicion}

\subsection{El experimento de Anfinsen y la paradoja de Levinthal}

La premisa de toda predicción es que la secuencia contiene la información necesaria para el plegamiento. \textcite{c15-anfinsen1973} lo demostró con la ribonucleasa A: al tratarla con urea y $\beta$-mercaptoetanol, la enzima se desplegaba y sus cuatro puentes disulfuro se rompían; al retirar lentamente los agentes desnaturalizantes, la proteína recuperaba espontáneamente su estructura y su actividad. De las 105 formas posibles de emparejar ocho cisteínas en cuatro puentes, la cadena elegía por sí sola la correcta. Anfinsen formuló así la \emph{hipótesis termodinámica}: la estructura nativa es la de mínima energía libre en las condiciones fisiológicas, y está determinada por la secuencia. Le valió el premio Nobel de Química de 1972.

Pero si la estructura nativa es un mínimo de energía, ¿cómo lo encuentra la proteína? Cyrus Levinthal señaló en 1969 una dificultad aritmética. Supongamos, de forma muy conservadora, que cada residuo puede adoptar solo tres conformaciones. Una proteína de 100 residuos tiene entonces
\begin{equation}\label{eq:15-levinthal}
\begin{aligned}
3^{100}&\approx 5{,}2\times10^{47}\ \text{conformaciones},\\
t_{\text{búsqueda}} &\approx 5{,}2\times10^{47}\times\SI{e-13}{\second}\approx 5\times10^{34}\ \si{\second}\approx 1{,}6\times10^{27}\ \text{años},
\end{aligned}
\end{equation}
si explorara una conformación cada \SI{0,1}{\pico\second}, el tiempo de una vibración molecular. Es un tiempo $10^{17}$ veces mayor que la edad del universo, y sin embargo las proteínas pequeñas se pliegan en milisegundos. Esta es la \emph{paradoja de Levinthal}.

\begin{simbolos}
\simbolo{$3^{100}$}{Número de conformaciones si cada uno de 100 residuos tiene 3 estados posibles.}
\simbolo{$\SI{e-13}{\second}$}{Tiempo optimista para visitar una conformación (una vibración de enlace).}
\simbolo{$t_{\text{búsqueda}}$}{Tiempo de una búsqueda exhaustiva y aleatoria.}
\end{simbolos}

La resolución es que el plegamiento no es una búsqueda aleatoria. \textcite{c15-zwanzig1992} mostraron con un modelo sencillo que basta un pequeño sesgo energético a favor de las conformaciones locales correctas para reducir el tiempo de búsqueda de escalas cósmicas a segundos. \textcite{c15-dill1997} lo expresaron con una imagen que se ha vuelto canónica: el paisaje de energía de una proteína no es un campo de golf plano con un único hoyo, sino un \emph{embudo} rugoso, en el que casi cualquier movimiento que reduzca la energía acerca a la cadena al estado nativo (\cref{fig:15-embudo}).

\begin{figure}[t]
\centering
\begin{tikzpicture}
\begin{groupplot}[group style={group size=2 by 1, horizontal sep=1.1cm},
   curso, width=0.52\linewidth, height=4.8cm, xmin=-1, xmax=1, ymin=-10, ymax=4,
   xtick=\empty, ytick=\empty, grid=none, xlabel={coordenada conformacional}]
\nextgroupplot[title={campo de golf (Levinthal)}, ylabel={energía libre}]
  \addplot[azul, very thick, domain=-1:1, samples=400] {0.25*sin(deg(40*x)) - 9*exp(-((x-0.45)/0.025)^2)};
  \node[etiqueta,anchor=north] at (axis cs:0.45,-9.3) {nativo};
  \node[etiqueta,align=center] at (axis cs:-0.35,2.3) {plano: ninguna pista\\sobre dónde está el hoyo};
\nextgroupplot[title={embudo (Dill y Chan)}]
  \addplot[naranja, very thick, domain=-1:1, samples=400] {-9 + 11*abs(x)^1.3 + 0.5*sin(deg(30*x))*abs(x)};
  \node[etiqueta,anchor=north] at (axis cs:0,-9.3) {nativo};
  \draw[flecha=tinta2] (axis cs:-0.85,1.8) -- (axis cs:-0.35,-4.2);
  \draw[flecha=tinta2] (axis cs:0.85,1.8) -- (axis cs:0.35,-4.2);
  \node[etiqueta,align=center] at (axis cs:0,2.6) {muchos caminos\\cuesta abajo};
\end{groupplot}
\end{tikzpicture}
\caption[Paisajes de energía del plegamiento]{Dos paisajes de energía libre idealizados, en una sola coordenada. \textbf{Izquierda}: si el paisaje fuera plano salvo por un hoyo estrecho, la cadena tendría que encontrarlo por búsqueda aleatoria, y tardaría más que la edad del universo (\cref{eq:15-levinthal}). \textbf{Derecha}: en un embudo, cada paso que reduce la energía acerca la cadena al estado nativo; la rugosidad produce mínimos locales transitorios, pero no impide llegar al fondo. La misma imagen de embudo reaparecerá en la \cref{sec:15-docking} para la unión de ligandos.}
\label{fig:15-embudo}
\end{figure}

Resolver la paradoja no resolvió el problema computacional. Simular físicamente el plegamiento desde la secuencia exige campos de fuerza exactos y tiempos de simulación inalcanzables para casi todas las proteínas. Durante cinco décadas, la predicción avanzó por otro camino: aprovechar la información evolutiva.

\subsection{Modelado por homología}

El ejemplo de las globinas mostró que proteínas con un \SI{28}{\percent} de identidad comparten el mismo plegamiento. Si nuestra proteína tiene un homólogo de estructura conocida (una \emph{plantilla}), podemos construir un modelo copiando su esqueleto. Es el \emph{modelado comparativo} o por homología, cuyo representante más influyente es MODELLER \parencite{c15-sali1993}. Su idea no es copiar coordenadas sino \emph{satisfacer restricciones espaciales}: del alineamiento con una o varias plantillas se derivan densidades de probabilidad para distancias y diedros del modelo; se combinan con términos estereoquímicos (longitudes y ángulos de enlace, contactos no permitidos) en una función objetivo, y se optimiza con gradientes conjugados y dinámica molecular con recocido simulado.

La calidad del modelo depende casi por completo del alineamiento y de la identidad con la plantilla. Por encima del \SI{50}{\percent} los modelos suelen ser muy buenos; entre el 30 y el \SI{50}{\percent} el núcleo es fiable pero lazos e inserciones no; por debajo del \SI{30}{\percent} los errores de alineamiento se vuelven la principal fuente de error. Y cuando no hay plantilla, el método simplemente no puede empezar.

\subsection{CASP: medir el progreso a ciegas}

Para saber qué métodos funcionan de verdad hace falta evaluarlos sin que los autores conozcan la respuesta. \textcite{c15-moult1995} organizaron en 1994 el primer CASP (\ingles{Critical Assessment of Structure Prediction}): los organizadores recogen secuencias de proteínas cuya estructura está a punto de resolverse experimentalmente, los grupos envían sus predicciones antes de que se publique, y evaluadores independientes las comparan con la estructura real. El experimento se repite cada dos años, y es el equivalente estructural del ensayo clínico doble ciego.

La medida principal de CASP es el GDT\_TS (\ingles{Global Distance Test}):
\begin{equation}\label{eq:15-gdt}
\mathrm{GDT\_TS}=\frac{100}{4}\left(P_1+P_2+P_4+P_8\right),
\end{equation}
donde $P_c$ es la mayor fracción de residuos del modelo que pueden superponerse a menos de $c$ \si{\angstrom} de sus posiciones reales. Como el TM-score, premia la fracción bien predicha e ignora los residuos muy desviados. Un GDT\_TS por encima de 90 se considera competitivo con una estructura experimental.

\begin{simbolos}
\simbolo{$P_c$}{Fracción máxima de C$_\alpha$ a menos de $c$ \si{\angstrom} de la referencia, optimizando la superposición para cada umbral.}
\simbolo{$c$}{Umbral de distancia: 1, 2, 4 y \SI{8}{\angstrom}.}
\end{simbolos}

\subsection{La señal de la coevolución}

La intuición de la apertura se formaliza así. Si en un alineamiento múltiple (\cref{cap:04}) las columnas $i$ y $j$ mutan de forma correlacionada, es probable que los residuos $i$ y $j$ estén en contacto, porque un cambio en uno exige un cambio compensatorio en el otro para mantener la estructura (\cref{fig:15-coevol}). La primera medida que viene a la mente es la \emph{información mutua} entre columnas,
\begin{equation}\label{eq:15-mi}
\mathrm{MI}_{ij}=\sum_{a,b} f_{ij}(a,b)\,\log\frac{f_{ij}(a,b)}{f_i(a)\,f_j(b)},
\end{equation}
donde $f_i(a)$ es la frecuencia del aminoácido $a$ en la columna $i$ y $f_{ij}(a,b)$ la frecuencia del par $(a,b)$ en las columnas $(i,j)$. Pero la información mutua confunde correlaciones directas e indirectas: si $i$ toca a $k$ y $k$ toca a $j$, las columnas $i$ y $j$ covarían aunque no se toquen, y además la filogenia (secuencias emparentadas que comparten mutaciones por herencia) añade correlaciones espurias.

\begin{simbolos}
\simbolo{$f_i(a)$}{Frecuencia (con pesos y pseudocuentas) del aminoácido $a$ en la columna $i$ del MSA.}
\simbolo{$f_{ij}(a,b)$}{Frecuencia conjunta del par $(a,b)$ en las columnas $i$ y $j$.}
\simbolo{$\mathrm{MI}_{ij}$}{Información mutua entre las columnas $i$ y $j$, en nats.}
\end{simbolos}

\begin{figure}[t]
\centering
% Coevolución en un MSA y contacto en la estructura
\begin{tikzpicture}[font=\sffamily\small,
  aa/.style={minimum width=4.6mm,minimum height=4.6mm,inner sep=0pt,font=\ttfamily\small,text=tinta}]
  \node[etiqueta,anchor=east] at (-0.1,0.0) {sec.\,1};
  \foreach \c/\ch in {0/M,1/K,2/V,3/L,4/A,5/E,6/G,7/R,8/D,9/K,10/T}{\node[aa] at (\c*0.5,0.0) {\ch};}
  \node[etiqueta,anchor=east] at (-0.1,-0.5) {sec.\,2};
  \foreach \c/\ch in {0/M,1/K,2/I,3/L,4/S,5/K,6/G,7/K,8/D,9/E,10/T}{\node[aa] at (\c*0.5,-0.5) {\ch};}
  \node[etiqueta,anchor=east] at (-0.1,-1.0) {sec.\,3};
  \foreach \c/\ch in {0/M,1/R,2/V,3/L,4/A,5/E,6/G,7/R,8/D,9/K,10/S}{\node[aa] at (\c*0.5,-1.0) {\ch};}
  \node[etiqueta,anchor=east] at (-0.1,-1.5) {sec.\,4};
  \foreach \c/\ch in {0/M,1/K,2/V,3/L,4/S,5/K,6/G,7/K,8/D,9/E,10/S}{\node[aa] at (\c*0.5,-1.5) {\ch};}
  \node[etiqueta,anchor=east] at (-0.1,-2.0) {sec.\,5};
  \foreach \c/\ch in {0/M,1/R,2/I,3/I,4/A,5/E,6/G,7/R,8/D,9/K,10/T}{\node[aa] at (\c*0.5,-2.0) {\ch};}
  \node[etiqueta,anchor=east] at (-0.1,-2.5) {sec.\,6};
  \foreach \c/\ch in {0/M,1/K,2/V,3/L,4/T,5/K,6/G,7/K,8/D,9/E,10/T}{\node[aa] at (\c*0.5,-2.5) {\ch};}
  \begin{scope}[on background layer]
    \fill[naranja!25,rounded corners=2pt] (2.28,0.3) rectangle (2.72,-2.8);
    \fill[aqua!25,rounded corners=2pt] (4.28,0.3) rectangle (4.72,-2.8);
  \end{scope}
  \node[etiqueta,text=naranja!70!black] at (2.5,0.6) {$i$};
  \node[etiqueta,text=aqua!60!black] at (4.5,0.6) {$j$};
  \draw[{Stealth}-{Stealth},tinta2] (2.5,-3.0) to[bend right=35] node[below,etiqueta]{E$\leftrightarrow$K compensados} (4.5,-3.0);
  \node[etiqueta,text width=5.3cm,anchor=north] at (2.5,-3.75) {cuando $i$ pasa de E a K, $j$ pasa de K a E: el puente salino se conserva con las cargas intercambiadas};
  % estructura plegada
  \begin{scope}[shift={(8.6,-1.2)}]
    \draw[tinta2,line width=2.2pt,rounded corners=8pt] (-2.2,1.5) -- (-1.2,1.7) -- (-0.2,0.9) -- (-1.0,-0.2) -- (0.2,-0.9) -- (1.2,0.1) -- (0.6,1.0) -- (1.6,1.6) -- (2.3,0.6) -- (1.8,-1.4);
    \node[circle,fill=naranja,minimum size=4mm,inner sep=0pt,font=\sffamily\bfseries\scriptsize,text=white] (ri) at (-0.2,0.9) {$i$};
    \node[circle,fill=aqua,minimum size=4mm,inner sep=0pt,font=\sffamily\bfseries\scriptsize,text=white] (rj) at (0.6,1.0) {$j$};
    \draw[rojo,thick,densely dashed] (ri) -- node[below,etiqueta,text=rojooscuro,yshift=-6pt,fill=white,inner sep=1pt]{contacto $<\SI{8}{\angstrom}$} (rj);
    \node[etiqueta,anchor=north] at (0,-1.7) {pliegue 3D: $i$ y $j$ se tocan};
  \end{scope}
  \draw[flecha=violeta] (5.7,-1.2) -- node[above,etiqueta,text=violeta]{DCA} node[below,etiqueta,text=violeta]{$J_{ij}$ grande} (6.6,-1.2);
\end{tikzpicture}

\caption[Coevolución y contactos]{La señal de la coevolución. \textbf{Izquierda}: en un alineamiento de seis homólogos, las columnas $i$ (naranja) y $j$ (aqua) cambian de forma coordinada: cuando $i$ es glutamato (E), $j$ es lisina (K), y viceversa, de modo que el puente salino entre ambos se conserva aunque los dos residuos cambien. \textbf{Derecha}: la explicación estructural más sencilla es que $i$ y $j$ están en contacto en el plegamiento. El análisis de acoplamiento directo (DCA) estima un acoplamiento $J_{ij}$ por par, separando la covariación directa de la transmitida a través de terceros residuos.}
\label{fig:15-coevol}
\end{figure}

\textcite{c15-morcos2011} resolvieron el problema con el \emph{análisis de acoplamiento directo} (DCA). En lugar de medir correlaciones par a par, ajustan a todo el MSA un modelo probabilístico global, un modelo de Potts de la física estadística,
\begin{equation}\label{eq:15-potts}
P(a_1,\dots,a_L)=\frac1Z\exp\!\Big(\sum_{i<j}J_{ij}(a_i,a_j)+\sum_i h_i(a_i)\Big),
\end{equation}
cuyos parámetros de acoplamiento $J_{ij}$ capturan solo las dependencias \emph{directas}, porque las indirectas ya quedan explicadas por las cadenas de acoplamientos intermedios. Con una aproximación de campo medio, los $J_{ij}$ se obtienen invirtiendo la matriz de covarianzas del alineamiento. Aplicado a más de un centenar de familias, los pares con mayor acoplamiento directo resultaron ser, en su gran mayoría, contactos reales, y con unas decenas de ellos como restricciones fue posible plegar proteínas sin plantilla.

\begin{simbolos}
\simbolo{$P(a_1,\dots,a_L)$}{Probabilidad de una secuencia de longitud $L$ bajo el modelo de la familia.}
\simbolo{$h_i(a)$}{Campo local: preferencia de la columna $i$ por el aminoácido $a$.}
\simbolo{$J_{ij}(a,b)$}{Acoplamiento directo entre el aminoácido $a$ en $i$ y el $b$ en $j$.}
\simbolo{$Z$}{Constante de normalización (función de partición).}
\end{simbolos}

\begin{cuidado}[La coevolución necesita muchas secuencias]
El número de parámetros de la \cref{eq:15-potts} crece como $L^2\times 20^2$: para una proteína de 200 residuos, unos ocho millones. Estimarlos requiere alineamientos con miles de secuencias \emph{efectivas} (tras reducir el peso de las muy parecidas). Para familias pequeñas, proteínas huérfanas o anticuerpos, la señal es débil, y esa limitación reaparece en todos los métodos que dependen del MSA, incluido AlphaFold2.
\end{cuidado}

\subsection{AlphaFold2}

En CASP14 (2020), un sistema de DeepMind obtuvo predicciones cuya exactitud, para la mayoría de las dianas, era comparable a la de las estructuras experimentales \parencite{c15-kryshtafovych2021}. \textcite{c15-jumper2021} describieron el método, AlphaFold2, e informaron una exactitud mediana del esqueleto de \SI{0,96}{\angstrom} (RMSD de los C$_\alpha$ sobre el \SI{95}{\percent} de los residuos), frente a \SI{2,8}{\angstrom} del siguiente mejor método. La diferencia no fue un avance incremental: fue un cambio de régimen.

AlphaFold2 no abandona la coevolución: la aprende. En lugar de extraer contactos con un modelo estadístico fijo y luego plegar, entrena una red neuronal de extremo a extremo que lee el MSA y produce directamente coordenadas atómicas. Su arquitectura tiene tres partes (\cref{fig:15-af2}).

\begin{figure}[t]
\centering
\resizebox{\linewidth}{!}{% Arquitectura estilizada de AlphaFold2
\begin{tikzpicture}[font=\sffamily\small,
  bloque/.style={caja=#1,minimum height=1.0cm,text width=2.25cm},
  bloque/.default=azul]
  % entrada
  \node[caja=gris,text width=1.7cm] (sec) at (0,0.55) {secuencia\\[-1pt]{\scriptsize\ttfamily MEEPQSD\ldots}};
  \node[caja=gris,text width=1.7cm,font=\sffamily\scriptsize] (bd) at (0,-1.55) {bases de datos\\(UniRef, BFD, PDB)};
  \node[bloque=aqua,text width=1.95cm] (msa) at (3.0,0.55) {MSA\\{\scriptsize $N_{\mathrm{sec}}\times N_{\mathrm{res}}\times c_m$}};
  \node[bloque=aqua,text width=1.95cm] (par) at (3.0,-1.55) {pares\\{\scriptsize $N_{\mathrm{res}}\times N_{\mathrm{res}}\times c_z$}};
  \draw[flecha] (sec) -- (msa);
  \draw[flecha] (sec) -- node[etiqueta,left]{búsqueda} (bd);
  \draw[flecha] (bd.north east) -- node[etiqueta,above,sloped,pos=0.6]{homólogos} (msa.south west);
  \draw[flecha] (bd) -- node[etiqueta,below]{plantillas} (par);
  % Evoformer
  \node[draw=azul!70!black,fill=azul!8,rounded corners=4pt,thick,minimum width=3.2cm,minimum height=3.75cm] (evo) at (6.45,-0.5) {};
  \node[font=\sffamily\bfseries\small,text=azulprofundo,anchor=north] at (evo.north) {Evoformer ($\times48$)};
  \node[caja=azul,font=\sffamily\scriptsize,text width=2.6cm] (e1) at (6.45,0.45) {atención por filas y columnas del MSA};
  \node[caja=azul,font=\sffamily\scriptsize,text width=2.6cm] (e2) at (6.45,-1.55) {actualizaciones triangulares de pares};
  \draw[flecha=azulprofundo] ([xshift=-4mm]e1.south) -- node[etiqueta,left]{producto} node[etiqueta,left,yshift=-2.5mm]{externo} ([xshift=-4mm]e2.north);
  \draw[flecha=azulprofundo] ([xshift=4mm]e2.north) -- node[etiqueta,right]{sesgo de} node[etiqueta,right,yshift=-2.5mm]{atención} ([xshift=4mm]e1.south);
  \draw[flecha] (msa.east) -- (evo.west |- msa.east);
  \draw[flecha] (par.east) -- (evo.west |- par.east);
  % módulo de estructura
  \node[bloque=naranja,text width=2.3cm] (est) at (10.9,-0.5) {módulo de estructura ($\times8$)\\{\scriptsize atención de punto invariante}};
  \draw[flecha] (evo.east |- est.west) ++(0,0.25) -- node[etiqueta,above]{individual} ([yshift=2.5mm]est.west);
  \draw[flecha] (evo.east |- est.west) ++(0,-0.25) -- node[etiqueta,below]{pares} ([yshift=-2.5mm]est.west);
  % salidas
  \node[caja=verde,text width=2.2cm,font=\sffamily\scriptsize] (xyz) at (10.9,1.55) {coordenadas 3D\\(marcos + torsiones)};
  \node[caja=amarillo,text width=2.2cm,font=\sffamily\scriptsize] (conf) at (10.9,-2.55) {confianza:\\pLDDT, PAE, pTM};
  \draw[flecha] (est) -- (xyz);
  \draw[flecha] (est) -- (conf);
  % reciclado
  \draw[flecha=violeta,rounded corners=6pt,dashed] (xyz.north) -- ++(0,0.45) -| node[pos=0.25,above,etiqueta,text=violeta]{reciclado (3 veces): salidas vuelven a entrar} (msa.north);
\end{tikzpicture}
}
\caption[Arquitectura de AlphaFold2]{Arquitectura simplificada de AlphaFold2 \parencite{c15-jumper2021}. La secuencia se usa para buscar homólogos (MSA) y plantillas. Dos representaciones se refinan juntas en 48 bloques Evoformer: la del MSA ($N_{\mathrm{sec}}\times N_{\mathrm{res}}$) y la de pares ($N_{\mathrm{res}}\times N_{\mathrm{res}}$); la información fluye del MSA a los pares mediante un producto externo, y de los pares al MSA como sesgo de la atención. El módulo de estructura convierte la representación individual y la de pares en un marco rígido (rotación y traslación) por residuo más sus ángulos de torsión, mediante la atención de punto invariante. Las salidas, incluida la estructura, se reciclan tres veces. Junto a las coordenadas, la red predice su propia confianza (pLDDT y PAE).}
\label{fig:15-af2}
\end{figure}

\paragraph{Entradas} La secuencia se busca en grandes bases de datos de secuencias con herramientas como las del \cref{cap:04} (\cmd{jackhmmer}, \cmd{HHblits}) para construir un MSA, y en el PDB para encontrar plantillas. El MSA se inicializa como un tensor $N_{\mathrm{sec}}\times N_{\mathrm{res}}\times c_m$ (una incrustación de dimensión $c_m$ por cada letra del alineamiento) y la representación de pares como un tensor $N_{\mathrm{res}}\times N_{\mathrm{res}}\times c_z$, en el que la celda $(i,j)$ acabará codificando todo lo que la red ``cree'' sobre la relación espacial entre los residuos $i$ y $j$.

\paragraph{Evoformer y atención} El corazón es la \emph{atención}, el mecanismo de los \ingles{transformers}. Para un conjunto de vectores de entrada $\mathbf x_1,\dots,\mathbf x_n$, cada uno se proyecta en una consulta $\mathbf q_i$, una clave $\mathbf k_j$ y un valor $\mathbf v_j$, y la salida en la posición $i$ es un promedio ponderado de los valores:
\begin{equation}\label{eq:15-atencion}
\mathbf o_i=\sum_{j=1}^{n}\alpha_{ij}\,\mathbf v_j,\qquad
\alpha_{ij}=\frac{\exp\!\big(\mathbf q_i^\top\mathbf k_j/\sqrt{c}+b_{ij}\big)}{\sum_{j'}\exp\!\big(\mathbf q_i^\top\mathbf k_{j'}/\sqrt{c}+b_{ij'}\big)}.
\end{equation}

\begin{simbolos}
\simbolo{$\mathbf q_i,\mathbf k_j,\mathbf v_j$}{Consulta, clave y valor: proyecciones lineales aprendidas de las entradas.}
\simbolo{$\alpha_{ij}$}{Peso de atención de la posición $i$ sobre la $j$; los pesos de cada $i$ suman 1 (\ingles{softmax}).}
\simbolo{$c$}{Dimensión de las consultas y claves; $\sqrt c$ estabiliza la escala del producto escalar.}
\simbolo{$b_{ij}$}{Sesgo aditivo; en AlphaFold2 procede de la representación de pares.}
\end{simbolos}

La atención permite que cada residuo ``consulte'' a todos los demás y decida a cuáles hacer caso. En el Evoformer se aplica en varias direcciones: por filas del MSA (cada secuencia mira a lo largo de sus residuos, con el sesgo $b_{ij}$ tomado de la representación de pares, de modo que la geometría supuesta guía la lectura del alineamiento), por columnas (cada posición compara lo que ocurre en distintas secuencias, que es donde vive la coevolución) y sobre la propia representación de pares. Para esta última, \textcite{c15-jumper2021} introdujeron las \emph{actualizaciones triangulares}: la celda $(i,j)$ se actualiza combinando las celdas $(i,k)$ y $(k,j)$ para todos los $k$. La motivación es geométrica: las distancias deben cumplir la desigualdad triangular, $d_{ij}\le d_{ik}+d_{kj}$, y hacer que cada par ``razone'' sobre los triángulos que forma con terceros residuos impone consistencia tridimensional antes de que exista ninguna coordenada. Finalmente, un \emph{producto externo} promediado sobre las secuencias del MSA transfiere la información de coevolución a la representación de pares.

\paragraph{Módulo de estructura} Cada residuo se representa como un cuerpo rígido, con un marco local (rotación y traslación, definidas por N, C$_\alpha$ y C) que inicialmente está en el origen. Durante ocho bloques con pesos compartidos, la \emph{atención de punto invariante} (IPA) actualiza esos marcos: es una atención en la que las consultas y claves incluyen puntos tridimensionales expresados en el marco de cada residuo, construida para que el resultado no cambie si se rota o traslada toda la estructura. Al final se predicen los ángulos de torsión de las cadenas laterales, y con ellos las coordenadas de todos los átomos pesados.

\paragraph{Entrenamiento y reciclado} La función de pérdida principal, FAPE (\ingles{frame aligned point error}), compara las posiciones atómicas predichas y reales expresadas en el marco local de cada residuo, lo que la hace sensible también a la quiralidad. La red se entrenó con estructuras del PDB y, en una segunda fase, con predicciones propias de alta confianza sobre secuencias sin estructura (autodestilación). La salida completa se vuelve a introducir como entrada tres veces (\ingles{recycling}), lo que permite refinar progresivamente la predicción.

\subsection{Confianza: pLDDT y PAE}

Una predicción sin estimación de su error sería peligrosa. AlphaFold2 aprende a predecir, además de la estructura, cuánto se equivoca.

\begin{definicion}{lDDT y pLDDT}{plddt}
El lDDT \parencite{c15-mariani2013} compara un modelo con la referencia \emph{sin superponerlos}. Para cada átomo $i$, se consideran los pares $(i,j)$ con distancia de referencia $d^{\mathrm{ref}}_{ij}<R_0=\SI{15}{\angstrom}$, y se cuenta la fracción de ellos cuya distancia en el modelo se conserva dentro de cada tolerancia:
\begin{equation}\label{eq:15-lddt}
\mathrm{lDDT}_i=\frac14\sum_{\tau\in\{0{,}5,\,1,\,2,\,4\}}\frac{\big|\{j:\ d^{\mathrm{ref}}_{ij}<R_0,\ |d^{\mathrm{mod}}_{ij}-d^{\mathrm{ref}}_{ij}|<\tau\}\big|}{\big|\{j:\ d^{\mathrm{ref}}_{ij}<R_0\}\big|}.
\end{equation}
El pLDDT es la predicción que hace la red de $100\times\mathrm{lDDT}_i$ (calculado sobre los C$_\alpha$) para cada residuo, sin conocer la estructura real.
\end{definicion}

\begin{simbolos}
\simbolo{$d^{\mathrm{ref}}_{ij},\ d^{\mathrm{mod}}_{ij}$}{Distancias entre los átomos $i$ y $j$ en la estructura de referencia y en el modelo.}
\simbolo{$R_0$}{Radio de inclusión: solo cuentan los vecinos a menos de \SI{15}{\angstrom} en la referencia.}
\simbolo{$\tau$}{Tolerancias de 0,5, 1, 2 y \SI{4}{\angstrom}, promediadas.}
\end{simbolos}

Por ser local y no requerir superposición, el lDDT evalúa bien proteínas con varios dominios que se mueven entre sí, y el pLDDT hereda esa propiedad: mide si el entorno local de cada residuo está bien predicho, no si los dominios están bien colocados unos respecto a otros. Por convención, se interpreta en cuatro bandas: muy alta ($>90$), segura (70--90), baja (50--70) y muy baja ($<50$). Las regiones con pLDDT muy bajo suelen corresponder a segmentos intrínsecamente desordenados, que no tienen una estructura única que predecir.

Para la disposición relativa de dominios, AlphaFold2 ofrece el \emph{error alineado predicho} (PAE). El elemento $\mathrm{PAE}_{ij}$ es el error esperado, en \si{\angstrom}, en la posición del residuo $j$ si se superponen la predicción y la estructura real sobre el marco local del residuo $i$. Una matriz de PAE con bloques de valores bajos en la diagonal indica dominios bien definidos; valores altos \emph{fuera} de los bloques indican que la posición relativa de esos dominios es incierta.

\begin{figure}[t]
\centering
\begin{tikzpicture}
\begin{axis}[capquince plddt]
  \fill[azulprofundo!22] (axis cs:1,90) rectangle (axis cs:393,100);
  \fill[azulclaro!40] (axis cs:1,70) rectangle (axis cs:393,90);
  \fill[amarillo!25] (axis cs:1,50) rectangle (axis cs:393,70);
  \fill[naranja!18] (axis cs:1,0) rectangle (axis cs:393,50);
  \addplot[tinta, thick] table[x=res,y=plddt] {figuras/cap15/plddt.dat};
  \node[etiqueta,anchor=west,text=azulprofundo] at (axis cs:395,95) {muy alta};
  \node[etiqueta,anchor=west,text=azulprofundo] at (axis cs:395,80) {segura};
  \node[etiqueta,anchor=west,text=amarillo!30!black] at (axis cs:395,60) {baja};
  \node[etiqueta,anchor=west,text=naranja!70!black] at (axis cs:395,25) {muy baja};
  \draw[decorate,decoration={brace,amplitude=4pt},violeta,thick] (axis cs:94,102) -- node[above=3pt,etiqueta,text=violeta]{dominio de unión a ADN (94--292)} (axis cs:292,102);
  \draw[decorate,decoration={brace,amplitude=4pt},violeta,thick] (axis cs:325,102) -- node[above=3pt,etiqueta,text=violeta]{tetramerización} (axis cs:356,102);
  \draw[decorate,decoration={brace,amplitude=4pt},gris,thick] (axis cs:1,102) -- node[above=3pt,etiqueta]{transactivación} (axis cs:90,102);
\end{axis}
\end{tikzpicture}
\caption[pLDDT de la p53 humana]{pLDDT por residuo del modelo de AlphaFold DB para la proteína supresora de tumores p53 humana (UniProt P04637, versión 6 del modelo), con las cuatro bandas de confianza sombreadas. El dominio de unión a ADN (pLDDT medio 95,3) y el de tetramerización (90,9) se predicen con confianza muy alta. La región N-terminal de transactivación (48,4), el enlace entre dominios (47,4) y el extremo C-terminal (43,4) caen en la banda muy baja: son regiones intrínsecamente desordenadas, que en la célula solo adoptan estructura al unirse a otras proteínas. Datos: AlphaFold DB, procesados con \archivo{figuras/cap15/generar.py}.}
\label{fig:15-plddt}
\end{figure}

\begin{ejemplo}{Leer un modelo de AlphaFold DB: p53}{p53}
La p53 humana (393 residuos) es un caso de estudio ideal porque mezcla dominios compactos con largas regiones desordenadas. En su modelo de AlphaFold DB, el \SI{52,7}{\percent} de los residuos tiene pLDDT muy alto, el \SI{7,1}{\percent} seguro, el \SI{10,4}{\percent} bajo y el \SI{29,8}{\percent} muy bajo (\cref{fig:15-plddt}).

La matriz de PAE (\cref{fig:15-pae}) añade la información que el pLDDT no da. Dentro del dominio de unión a ADN, el PAE medio es de \SI{3,6}{\angstrom}, y dentro del de tetramerización, de \SI{2,9}{\angstrom}; pero el PAE medio de las posiciones del dominio de tetramerización vistas desde el marco del dominio de unión a ADN es de \SI{20,1}{\angstrom}. Traducción: cada dominio está bien predicho, pero su orientación relativa en el modelo es arbitraria, y sería un error medir distancias entre ellos.

¿Es fiable el dominio de alta confianza? Comparando sus C$_\alpha$ (residuos 94--292) con la estructura cristalográfica \texttt{2OCJ}, 194 pares tienen un RMSD de \SI{0,51}{\angstrom} y un TM-score de 0,99. Es una concordancia casi perfecta, aunque conviene una advertencia honesta: esa estructura, depositada en 2007, pudo formar parte del conjunto de entrenamiento. La evaluación justa de un predictor exige estructuras publicadas después de su entrenamiento, que es exactamente lo que garantiza CASP.
\end{ejemplo}

\begin{figure}[t]
\centering
\includegraphics[width=0.62\linewidth]{cap15/pae_p53.pdf}
\caption[Matriz de PAE de la p53 humana]{Matriz de error alineado predicho (PAE) del modelo de AlphaFold DB para la p53 humana. La celda $(i,j)$ es el error esperado en la posición del residuo $j$ tras superponer sobre el residuo $i$; los tonos oscuros indican confianza. Los dos bloques oscuros de la diagonal (recuadros) son el dominio de unión a ADN (DBD) y el de tetramerización (TET): cada uno es rígido y está bien definido. Las regiones claras fuera de ellos indican que la posición relativa de los dominios, y la de las regiones desordenadas, es desconocida. Datos: AlphaFold DB (\texttt{AF-P04637-F1}, v6).}
\label{fig:15-pae}
\end{figure}

\begin{codigo}[alphafold\_db.py]
import requests

acc = "P04637"                                   # p53 humana
api = f"https://alphafold.ebi.ac.uk/api/prediction/{acc}"
meta = requests.get(api, timeout=30).json()[0]
pdb = requests.get(meta["pdbUrl"], timeout=60).text
plddt = {}
for linea in pdb.splitlines():                   # pLDDT = columna B
    if linea.startswith("ATOM") and linea[12:16] == " CA ":
        plddt[int(linea[22:26])] = float(linea[60:66])
bajos = [r for r, v in plddt.items() if v < 50]
print(f"{len(plddt)} residuos; {len(bajos)} con pLDDT < 50")
\end{codigo}

\begin{cuidado}[pLDDT alto no significa estructura correcta en la célula]
El pLDDT mide la confianza de la red en una conformación, no la verdad biológica. Un modelo puede tener pLDDT alto y representar solo uno de varios estados funcionales (activo o inactivo, con o sin ligando), porque AlphaFold2 tiende a predecir la conformación más representada en su entrenamiento. Tampoco modela ligandos, iones, modificaciones postraduccionales ni otras cadenas. Y un pLDDT bajo no es un fracaso: a menudo es la mejor predicción posible, la de que la región está desordenada.
\end{cuidado}

\subsection{AlphaFold DB, RoseTTAFold, ESMFold, AlphaFold3 y ColabFold}

\paragraph{AlphaFold DB} Con el método publicado, DeepMind y el EMBL-EBI calcularon modelos para proteomas completos y los ofrecieron gratuitamente. \textcite{c15-varadi2022} describieron la primera versión de la AlphaFold Protein Structure Database, con los proteomas de referencia de humanos y de decenas de organismos modelo, y la ampliación posterior a prácticamente todo UniProt, con más de 200 millones de modelos, cambió la escala del campo: por primera vez, la mayoría de las proteínas conocidas tiene un modelo estructural. Como vimos, el PDB integra hoy estos modelos calculados en su portal, separados de las estructuras experimentales \parencite{c15-burley2023}.

\paragraph{RoseTTAFold} Casi simultáneamente, \textcite{c15-baek2021} presentaron RoseTTAFold, una red de ``tres pistas'' en la que la información de la secuencia (1D), de las distancias entre residuos (2D) y de las coordenadas (3D) fluye y se actualiza en paralelo. Con una exactitud cercana a la de AlphaFold2 y un código abierto desde el principio, mostró además que la misma arquitectura podía predecir complejos de proteínas.

\paragraph{ESMFold} Tanto AlphaFold2 como RoseTTAFold dependen de un MSA, cuya construcción es lenta y que falla para proteínas con pocos homólogos. \textcite{c15-lin2023} sustituyeron el MSA por un \emph{modelo de lenguaje de proteínas}, ESM-2, un \ingles{transformer} con hasta 15\,000 millones de parámetros entrenado únicamente a predecir aminoácidos enmascarados en millones de secuencias. Para acertar qué aminoácido falta, el modelo tiene que aprender implícitamente qué residuos interactúan, es decir, la coevolución, y esa información queda en sus representaciones internas. ESMFold acopla ESM-2 a un módulo de plegamiento y predice estructuras a partir de una sola secuencia, hasta un orden de magnitud más rápido que AlphaFold2, a cambio de una exactitud algo menor en proteínas con pocos homólogos. Su velocidad permitió plegar más de 600 millones de proteínas metagenómicas.

\paragraph{AlphaFold3} \textcite{c15-abramson2024} extendieron el enfoque a casi toda la química de la célula: proteínas, ácidos nucleicos, ligandos pequeños, iones y residuos modificados, en un mismo complejo. Los cambios de arquitectura son profundos: el Evoformer se reemplaza por un módulo más ligero (Pairformer) que procesa el MSA de forma reducida, y el módulo de estructura se reemplaza por un \emph{módulo de difusión} que genera directamente coordenadas atómicas partiendo de ruido y eliminándolo paso a paso, sin marcos ni torsiones. En las pruebas de sus autores, AlphaFold3 superó a los métodos de \ingles{docking} clásicos en la predicción de complejos proteína-ligando, un resultado que conecta directamente con la \cref{sec:15-docking}.

\paragraph{ColabFold} \textcite{c15-mirdita2022} hicieron AlphaFold2 accesible a cualquier investigador: reemplazaron la búsqueda de homólogos por MMseqs2 en un servidor, lo que acelera la construcción del MSA en uno o dos órdenes de magnitud, y lo empaquetaron en cuadernos de Google Colab y en una herramienta de línea de órdenes. Es la forma más sencilla de predecir la estructura de una proteína o de un complejo sin infraestructura propia.

\begin{consola}[Predicción con ColabFold en local]
# instalación: https://github.com/sokrypton/ColabFold
colabfold_batch --num-recycle 3 --amber --use-gpu-relax \
    p53_dbd.fasta resultados/
ls resultados/
# p53_dbd_unrelaxed_rank_001_..._model_3_seed_000.pdb
# p53_dbd_scores_rank_001_..._model_3_seed_000.json  (pLDDT, PAE)
# p53_dbd_coverage.png  p53_dbd_pae.png  p53_dbd_plddt.png
\end{consola}

\begin{table}[t]
\centering\small
\caption[Medidas de confianza de los predictores]{Medidas de confianza que acompañan a los modelos de AlphaFold y afines, y la pregunta que responde cada una.}
\label{tab:15-confianza}
\begin{tabularx}{\linewidth}{@{}>{\raggedright\arraybackslash}p{1.6cm}>{\raggedright\arraybackslash}p{2.6cm}>{\raggedright\arraybackslash}X@{}}
\toprule
\textbf{Medida} & \textbf{Alcance} & \textbf{Pregunta que responde}\\
\midrule
pLDDT & por residuo (0--100) & ¿está bien predicho el entorno local de este residuo?\\
PAE & por par de residuos (\si{\angstrom}) & ¿está bien colocado $j$ respecto a $i$? Revela dominios y su orientación relativa.\\
pTM & global (0--1) & TM-score esperado del modelo completo.\\
ipTM & interfaz (0--1) & en complejos, ¿está bien predicha la interfaz entre cadenas?\\
\bottomrule
\end{tabularx}
\end{table}

\begin{historia}[El premio del problema de cincuenta años]
El problema del plegamiento se formuló, en su versión moderna, con el experimento de Anfinsen a principios de los sesenta. CASP se creó en 1994 precisamente porque los avances anunciados rara vez resistían una evaluación a ciegas. Veintiséis años después, en CASP14, los organizadores declararon que, para las proteínas individuales, el problema estaba esencialmente resuelto \parencite{c15-kryshtafovych2021}. En 2024, Demis Hassabis y John Jumper compartieron el premio Nobel de Química por AlphaFold, junto con David Baker, del grupo de RoseTTAFold, por el diseño computacional de proteínas.
\end{historia}



% ---------------------------------------------------------------------
\section{Acoplamiento molecular (docking)}\label{sec:15-docking}
% ---------------------------------------------------------------------
\enlacerepo{15.3 · Docking molecular básico}

\begin{intuicion}[Una llave en la oscuridad]
Imagine que debe encontrar, a oscuras, cuál de mil llaves abre una cerradura, y en qué posición entra. Puede palpar la forma de la cerradura, pero no ver. Para cada llave tiene que hacer dos cosas: \emph{probar} posiciones (girarla, inclinarla, empujarla, doblar las llaves articuladas) y \emph{juzgar} cada posición (¿encaja sin forzar?, ¿hace contacto con los pernos?). Si prueba pocas posiciones, puede descartar la llave correcta porque nunca la metió bien; si su tacto es torpe, puede aceptar una llave que entra pero no gira. El acoplamiento molecular tiene exactamente esos dos componentes, la búsqueda y la puntuación, y casi todos sus fracasos provienen de uno de ellos.
\end{intuicion}

\subsection{El problema}

El \emph{acoplamiento molecular} (\ingles{docking}) predice la posición, la orientación y la conformación (juntas, la \emph{pose}) de una molécula pequeña, el \emph{ligando}, unida a una macromolécula, el \emph{receptor}, y estima la fuerza de esa unión. Es la herramienta central del diseño de fármacos basado en estructura: permite cribar virtualmente millones de compuestos contra una diana, proponer cómo se une un inhibidor conocido o sugerir modificaciones químicas que mejoren su afinidad.

\textcite{c15-kuntz1982} formularon el primer método: representaron el bolsillo del receptor como un conjunto de esferas que rellenan su hueco y buscaron cómo superponer los átomos del ligando sobre los centros de esas esferas, evaluando la complementariedad de forma. Su programa, DOCK, estableció el principio que sigue vigente: el ligando debe encajar geométricamente (complementariedad de \emph{forma}) y sus grupos deben enfrentarse a los grupos adecuados del receptor (complementariedad \emph{química}: donadores frente a aceptores de puentes de hidrógeno, cargas opuestas, superficies hidrofóbicas juntas).

\begin{figure}[t]
\centering
\resizebox{\linewidth}{!}{% Esquema del docking: grados de libertad y ciclo búsqueda-puntuación
\begin{tikzpicture}[font=\sffamily\small]
  % receptor con bolsillo
  \fill[azulpalido,draw=azul!70!black,thick] plot[smooth cycle,tension=0.7] coordinates
    {(-0.2,0) (0.3,2.2) (1.6,3.0) (2.3,2.1) (2.8,1.3) (3.6,2.2) (4.7,2.5) (5.1,1.0) (4.6,-0.9) (2.4,-1.4) (0.6,-1.1)};
  \node[etiqueta,text=azulprofundo,font=\sffamily\bfseries\small] at (1.6,0.1) {receptor};
  \fill[amarillo!30] (3.2,1.55) circle (0.55);
  \node[etiqueta,text=amarillo!30!black] at (3.2,0.85) {bolsillo};
  % caja de búsqueda
  \draw[violeta,dashed,thick] (2.3,0.6) rectangle (4.3,2.9);
  \node[etiqueta,text=violeta,anchor=south] at (3.3,2.95) {caja de búsqueda};
  % ligando fuera
  \begin{scope}[shift={(7.2,2.6)}]
    \draw[tinta,line width=1.3pt] (0,0) -- (0.5,0.3) -- (1.0,0) -- (1.5,0.3);
    \draw[tinta,line width=1.3pt] (0.5,0.3) -- (0.5,0.85);
    \foreach \p/\c in {{(0,0)}/gris,{(0.5,0.3)}/gris,{(1.0,0)}/azul,{(1.5,0.3)}/rojo,{(0.5,0.85)}/rojo}{\fill[\c] \p circle (0.11);}
    \draw[-{Stealth[length=1.6mm]},naranja,thick] (1.25,-0.05) arc[start angle=-90,end angle=200,x radius=0.12,y radius=0.26];
    \node[etiqueta,text=naranja!80!black,anchor=west] at (1.5,-0.25) {torsión};
    \node[etiqueta,anchor=north] at (0.75,-0.4) {ligando flexible};
  \end{scope}
  \draw[flecha=tinta2] (7.0,2.45) to[bend left=15] node[below,etiqueta,pos=0.45,yshift=-2pt]{traslación + rotación} (3.9,1.6);
  \node[etiqueta,text width=4.2cm,anchor=north west,align=left] at (6.4,1.2) {grados de libertad:\\$3$ (posición) $+\,3$ (orientación)\\$+\,N_{\mathrm{rot}}$ (torsiones)};
  % ciclo
  \node[caja=violeta,text width=2.6cm] (bus) at (1.0,-2.6) {búsqueda\\{\scriptsize (Monte Carlo, GA, BFGS)}};
  \node[caja=naranja,text width=2.6cm] (pun) at (5.4,-2.6) {puntuación\\{\scriptsize $\Delta G$ aproximada}};
  \draw[flecha=tinta2] (bus.north east) to[bend left=18] node[above,etiqueta]{pose candidata} (pun.north west);
  \draw[flecha=tinta2] (pun.south west) to[bend left=18] node[below,etiqueta]{energía} (bus.south east);
  \node[caja=verde,text width=2.7cm] (sal) at (9.2,-2.6) {poses ordenadas\\{\scriptsize (kcal/mol)}};
  \draw[flecha=verde] (pun) -- (sal);
\end{tikzpicture}
}
\caption[Esquema del acoplamiento molecular]{Los dos componentes del \ingles{docking}. \textbf{Arriba}: el ligando flexible debe colocarse dentro de una caja de búsqueda que rodea el bolsillo del receptor; su pose se describe con tres coordenadas de posición, tres de orientación y un ángulo por cada enlace rotable, $6+N_{\mathrm{rot}}$ grados de libertad en total. \textbf{Abajo}: el algoritmo de búsqueda propone poses y la función de puntuación las evalúa; el ciclo se repite hasta converger, y la salida es una lista de poses ordenadas por su energía estimada.}
\label{fig:15-dockesq}
\end{figure}

Formalmente, si $\mathbf x=(\mathbf t,\mathbf r,\boldsymbol\tau)$ reúne la traslación $\mathbf t\in\R^3$, la orientación $\mathbf r$ (tres parámetros, por ejemplo un cuaternión unitario o un vector de rotación) y los $N_{\mathrm{rot}}$ ángulos de torsión $\boldsymbol\tau$ del ligando, el \ingles{docking} busca
\begin{equation}\label{eq:15-docking}
\mathbf x^\ast=\argmin_{\mathbf x\in\mathcal B\times SO(3)\times\mathbb T^{N_{\mathrm{rot}}}} S(\mathbf x),
\end{equation}
donde $S$ es una función de puntuación que aproxima la energía libre de unión y $\mathcal B$ es la caja de búsqueda. El receptor suele tratarse como rígido.

\begin{simbolos}
\simbolo{$\mathbf x$}{Pose del ligando: posición, orientación y torsiones.}
\simbolo{$\mathcal B$}{Caja de búsqueda (región del receptor donde se permite colocar el centro del ligando).}
\simbolo{$\mathbb T^{N_{\mathrm{rot}}}$}{Toro de dimensión $N_{\mathrm{rot}}$: cada torsión es un ángulo periódico.}
\simbolo{$S(\mathbf x)$}{Puntuación de la pose; más negativa significa unión más favorable.}
\end{simbolos}

Ambos factores son difíciles. El espacio de búsqueda de un ligando típico, con 5 a 10 enlaces rotables, tiene 11 a 16 dimensiones continuas y está lleno de mínimos locales. Y la energía libre de unión real incluye efectos que ninguna función rápida captura bien: la desolvatación de ligando y bolsillo, la pérdida de entropía conformacional, el papel de moléculas de agua concretas, la polarización.

\subsection{Encontrar el sitio de unión}

Antes de acoplar hay que saber dónde. Si existe una estructura del receptor con un ligando cocristalizado, el bolsillo es conocido y la caja se centra en él. Si no, los detectores de bolsillos buscan cavidades en la superficie con criterios geométricos (regiones cóncavas donde cabe una esfera de sonda pero que están rodeadas de proteína en muchas direcciones), a menudo combinados con la conservación evolutiva de los residuos que las tapizan: un bolsillo conservado en toda una familia probablemente tiene función. El docking ``ciego'', con una caja que abarca toda la proteína, es posible pero multiplica el espacio de búsqueda y las poses falsas.

\subsection{Funciones de puntuación}

Hay tres grandes familias de funciones de puntuación: las basadas en \emph{campos de fuerza} (términos de van der Waals y electrostáticos de la mecánica molecular), las \emph{empíricas} (suma de términos físicamente motivados cuyos pesos se ajustan por regresión a afinidades experimentales) y las \emph{basadas en conocimiento} (potenciales estadísticos derivados de las frecuencias de contactos observadas en complejos del PDB). La función de AutoDock Vina es un híbrido de las dos últimas y merece un estudio detallado, porque es la más usada en la práctica.

\textcite{c15-trott2010} definen la energía de interacción como una suma sobre pares de átomos $(i,j)$ del ligando y del receptor, separados por más de tres enlaces y a menos de \SI{8}{\angstrom}, de una función de la \emph{distancia de superficie}
\begin{equation}\label{eq:15-dsup}
d_{ij}=r_{ij}-R_{t_i}-R_{t_j},
\end{equation}
donde $r_{ij}$ es la distancia entre centros y $R_t$ el radio de van der Waals del tipo de átomo $t$ (por ejemplo, \SI{1,9}{\angstrom} para el carbono, 1,8 para el nitrógeno, 1,7 para el oxígeno). Los términos son
\begin{equation}\label{eq:15-vina}
\begin{aligned}
\mathrm{gauss}_1(d)&=e^{-(d/0{,}5)^2}, &\qquad
\mathrm{gauss}_2(d)&=e^{-((d-3)/2)^2},\\[2pt]
\mathrm{rep}(d)&=\begin{cases}d^2 & d<0\\0 & d\ge0\end{cases}, &
\mathrm{hidrof}(d)&=\begin{cases}1 & d<0{,}5\\ 1{,}5-d & 0{,}5\le d\le1{,}5\\ 0 & d>1{,}5\end{cases},\\[2pt]
\mathrm{puenteH}(d)&=\begin{cases}1 & d<-0{,}7\\ -d/0{,}7 & -0{,}7\le d\le0\\ 0 & d>0\end{cases},
\end{aligned}
\end{equation}
con pesos $w_{\mathrm{g1}}=-0{,}0356$, $w_{\mathrm{g2}}=-0{,}00516$, $w_{\mathrm{rep}}=0{,}840$, $w_{\mathrm{hidrof}}=-0{,}0351$ y $w_{\mathrm{H}}=-0{,}587$. El término hidrofóbico solo se aplica entre átomos hidrofóbicos, y el de puente de hidrógeno solo entre un donador y un aceptor. La energía intermolecular total se corrige por la flexibilidad del ligando:
\begin{equation}\label{eq:15-vinarot}
S=\frac{\sum_{i<j} \sum_{m} w_m\,h_m(d_{ij})}{1+w_{\mathrm{rot}}\,N_{\mathrm{rot}}},\qquad w_{\mathrm{rot}}=0{,}0585.
\end{equation}

\begin{simbolos}
\simbolo{$r_{ij}$}{Distancia entre los centros de los átomos $i$ (ligando) y $j$ (receptor).}
\simbolo{$R_t$}{Radio de van der Waals del tipo de átomo $t$.}
\simbolo{$d_{ij}$}{Distancia entre las superficies atómicas; negativa si las esferas se interpenetran.}
\simbolo{$h_m,\ w_m$}{Término $m$ de la \cref{eq:15-vina} y su peso ajustado.}
\simbolo{$N_{\mathrm{rot}}$}{Número de enlaces rotables del ligando (entre átomos pesados, no terminales).}
\simbolo{$S$}{Puntuación final, que Vina expresa en kcal/mol como estimación de $\Delta G$ de unión.}
\end{simbolos}

\begin{figure}[t]
\centering
\begin{tikzpicture}
\begin{axis}[curso, width=0.86\linewidth, height=5.6cm, xmin=2.4, xmax=8, ymin=-0.26, ymax=0.2,
   xlabel={distancia entre centros $r$ (\si{\angstrom})}, ylabel={energía del par (kcal/mol)},
   legend pos=south east, ytick={-0.2,-0.1,0,0.1,0.2}]
  \addplot[azul, very thick] table[x=r,y=ecc] {figuras/cap15/vina.dat};
  \addlegendentry{C$\cdots$C hidrofóbico}
  \addplot[naranja, very thick] table[x=r,y=ehb] {figuras/cap15/vina.dat};
  \addlegendentry{N$\cdots$O donador-aceptor}
  \addplot[tinta2, dashed, thin] coordinates {(2.4,0)(8,0)};
  \addplot[only marks, mark=*, mark size=1.8pt, azul] coordinates {(3.80,-0.071)};
  \addplot[only marks, mark=*, mark size=1.8pt, naranja] coordinates {(3.03,-0.224)};
  \node[etiqueta,anchor=north,text=azulprofundo] at (axis cs:3.9,-0.08) {mínimo \SI{3,8}{\angstrom}};
  \node[etiqueta,anchor=west,text=naranja!70!black] at (axis cs:3.1,-0.225) {mínimo \SI{3,0}{\angstrom}};
  \node[etiqueta,anchor=west] at (axis cs:2.45,0.15) {repulsión estérica};
\end{axis}
\end{tikzpicture}
\caption[Términos de la función de Vina]{Energía de un par de átomos según la función de puntuación de AutoDock Vina (\cref{eq:15-vina}), en función de la distancia entre centros. Para dos carbonos hidrofóbicos (azul), el mínimo, de $-0{,}071$ kcal/mol, está en \SI{3,8}{\angstrom}, justo donde se tocan sus superficies de van der Waals. Un par donador-aceptor de puente de hidrógeno (naranja) es tres veces más favorable y su mínimo se desplaza a \SI{3,0}{\angstrom}, la distancia típica N$\cdots$O de un puente de hidrógeno. A distancias cortas domina la repulsión cuadrática. Un ligando con decenas de contactos suma estas pequeñas contribuciones hasta alcanzar varias kcal/mol. Curvas calculadas con \archivo{figuras/cap15/generar.py}.}
\label{fig:15-vina}
\end{figure}

La \cref{fig:15-vina} muestra la forma de estos términos. Tienen una lógica clara: la atracción gaussiana estrecha premia el contacto íntimo; la ancha, la proximidad general; la repulsión castiga los choques; y los dos términos lineales recompensan los contactos hidrofóbicos y los puentes de hidrógeno. Llama la atención lo que \emph{falta}: no hay término electrostático explícito ni desolvatación. Los pesos se ajustaron para reproducir afinidades experimentales, así que esos efectos quedan absorbidos, de forma promedio, en los términos existentes. Es una función deliberadamente simple, y su virtud es la velocidad.

\begin{ejemplo}{De la puntuación a la constante de disociación}{vina}
Suponga que la suma intermolecular de una pose es de $-10{,}5$ kcal/mol y que el ligando tiene $N_{\mathrm{rot}}=5$ enlaces rotables. La \cref{eq:15-vinarot} divide por $1+0{,}0585\times5=1{,}29$, de modo que $S=-8{,}12$ kcal/mol: la flexibilidad le cuesta al ligando más de 2 kcal/mol, porque al unirse pierde la libertad de girar esos enlaces. Si interpretamos $S$ como energía libre de unión, la relación termodinámica $\Delta G=RT\ln K_d$ a \SI{298}{\kelvin} ($RT=0{,}593$ kcal/mol) da
\[
K_d = e^{\Delta G/RT}=e^{-8{,}12/0{,}593}\approx1{,}1\times10^{-6}\ \mathrm{M}\approx\SI{1,1}{\micro M}.
\]
Cada $RT\ln10=1{,}36$ kcal/mol adicionales multiplica la afinidad por diez: $-9$ kcal/mol corresponde a unos \SI{0,25}{\micro M} y $-12$ kcal/mol a unos \SI{1,6}{\nano M}. Pero no tome estas cifras al pie de la letra: el error típico de las funciones de puntuación al predecir afinidades absolutas es de unas 2 kcal/mol, es decir, más de un orden de magnitud en $K_d$. El \ingles{docking} ordena poses de un mismo ligando mucho mejor que afinidades entre ligandos distintos.
\end{ejemplo}

\subsection{Búsqueda conformacional}

Los algoritmos de búsqueda también son variados: búsqueda sistemática sobre una rejilla de ángulos, construcción incremental del ligando fragmento a fragmento dentro del bolsillo, algoritmos genéticos (el de AutoDock 4) y métodos de Monte Carlo. Vina usa una \emph{búsqueda local iterada}: parte de una pose aleatoria, la perturba (un desplazamiento, un giro o el cambio de una torsión), minimiza localmente la energía con el método cuasi-Newton BFGS usando las derivadas analíticas de la función de puntuación, y acepta o rechaza la nueva pose con el criterio de Metropolis,
\begin{equation}\label{eq:15-metropolis}
\Prob(\text{aceptar})=\min\!\big(1,\ e^{-(S_{\text{nueva}}-S_{\text{actual}})/T}\big),
\end{equation}
que acepta siempre las mejoras y ocasionalmente los empeoramientos, para escapar de mínimos locales. Varias de estas cadenas corren en paralelo (el parámetro \texttt{exhaustiveness} controla cuántas), y para acelerar la evaluación, la contribución de cada tipo de átomo del ligando se precalcula en una rejilla tridimensional sobre la caja de búsqueda. Con estas decisiones, \textcite{c15-trott2010} informaron una aceleración de unos dos órdenes de magnitud respecto a AutoDock 4 junto con una mayor exactitud en la predicción de poses.

\begin{simbolos}
\simbolo{$S_{\text{nueva}},\ S_{\text{actual}}$}{Puntuaciones de la pose propuesta (ya minimizada) y de la actual.}
\simbolo{$T$}{``Temperatura'' del algoritmo: controla cuánto se toleran los empeoramientos.}
\end{simbolos}

La versión 1.2 \parencite{c15-eberhardt2021} añadió enlaces de Python, la función de puntuación de AutoDock 4 como alternativa, el acoplamiento simultáneo de varios ligandos, el tratamiento de macrociclos flexibles y un protocolo con moléculas de agua explícitas. El flujo completo en Python es breve:

\begin{codigo}[vina\_docking.py]
from vina import Vina                           # AutoDock Vina >= 1.2

v = Vina(sf_name="vina", seed=42)
v.set_receptor("receptor.pdbqt")                # con hidrógenos polares
v.set_ligand_from_file("ligando.pdbqt")         # torsiones definidas
v.compute_vina_maps(center=[15.2, 53.9, 16.9],  # centro del bolsillo
                    box_size=[20, 20, 20])      # caja en angstroms
v.dock(exhaustiveness=16, n_poses=9)
print(v.energies(n_poses=5)[:, 0])              # kcal/mol por pose
v.write_poses("poses.pdbqt", n_poses=5, overwrite=True)
\end{codigo}

La preparación de los archivos \texttt{.pdbqt} (añadir hidrógenos, asignar tipos de átomo y cargas, definir las torsiones rotables del ligando) se hace con herramientas como Meeko u Open Babel, y es una fuente frecuente de errores silenciosos: un estado de protonación equivocado de una histidina del bolsillo o de un grupo ácido del ligando puede cambiar por completo el resultado.

\subsection{El embudo de unión y la validación por \ingles{redocking}}

La imagen del embudo de la \cref{fig:15-embudo} vuelve a ser útil. Si la función de puntuación es buena, las poses se ordenan como un embudo: cuanto más se parecen a la pose real, mejor puntuación tienen, y la mejor puntuada está cerca de la real. Si la función tiene un mínimo falso, un ``embudo trampa'' en otra región del bolsillo, la búsqueda puede converger allí con total confianza.

\begin{figure}[t]
\centering
\begin{tikzpicture}
\begin{groupplot}[group style={group size=2 by 1, horizontal sep=1.6cm},
   curso, height=5.4cm]
\nextgroupplot[width=0.47\linewidth, view={0}{90}, xmin=-3, xmax=3, ymin=-3, ymax=3,
   xlabel={coordenada de pose 1}, ylabel={coordenada de pose 2}, grid=none,
   colormap name=curso, point meta min=-9.5, point meta max=1, title={paisaje de puntuación},
   xtick=\empty, ytick=\empty, axis on top]
  \addplot3[surf, shader=interp, colormap={inv}{indices of colormap={4,3,2,1,0 of curso}}] table {figuras/cap15/paisaje.dat};
  \node[circle,draw=white,fill=naranja,inner sep=1.6pt] at (axis cs:0.6,0.4) {};
  \node[etiqueta,text=white,anchor=south] at (axis cs:0.6,0.55) {pose nativa};
  \node[circle,draw=white,fill=magenta,inner sep=1.6pt] at (axis cs:-1.7,-1.4) {};
  \node[etiqueta,text=white,anchor=north] at (axis cs:-1.7,-1.6) {mínimo falso};
\nextgroupplot[width=0.5\linewidth, xmin=0, xmax=12, ymin=-10.5, ymax=-3.5,
   xlabel={RMSD a la pose cristalográfica (\si{\angstrom})}, ylabel={puntuación (kcal/mol)},
   title={puntuación frente a RMSD}]
  \fill[verde!12] (axis cs:0,-10.5) rectangle (axis cs:2,-3.5);
  \addplot[only marks, mark=*, mark size=0.9pt, azul, fill opacity=0.5, draw opacity=0.5] table[x=rmsd,y=score] {figuras/cap15/poses.dat};
  \addplot[tinta2, dashed] coordinates {(2,-10.5)(2,-3.5)};
  \node[etiqueta,anchor=north west,text=verde!40!black] at (axis cs:0.1,-3.6) {$<\SI{2}{\angstrom}$};
  \node[etiqueta,anchor=north,text=magenta!70!black] at (axis cs:7.5,-9.1) {embudo trampa};
\end{groupplot}
\end{tikzpicture}
\caption[Embudo de docking]{El embudo de unión en un ejemplo simulado. \textbf{Izquierda}: paisaje de puntuación sobre dos coordenadas de pose (tonos oscuros = puntuación más favorable), con el mínimo global en la pose nativa y un mínimo local profundo en otra región del bolsillo. \textbf{Derecha}: 400 poses generadas por la búsqueda, con su puntuación frente a su RMSD respecto a la pose real. La mayoría forma un embudo cuyo fondo está a menos de \SI{2}{\angstrom} (franja verde): en este caso la mejor pose tiene un RMSD de \SI{0,53}{\angstrom} y 7 de las 10 mejores son correctas. El grupo en torno a \SI{7,5}{\angstrom} es un embudo trampa: si la función de puntuación lo favoreciera un poco más, el \ingles{docking} fallaría con una puntuación excelente. Datos simulados con \archivo{figuras/cap15/generar.py}.}
\label{fig:15-funnel}
\end{figure}

Por eso, antes de usar un protocolo de \ingles{docking} con compuestos nuevos, se \emph{valida} con un caso de respuesta conocida: el \ingles{redocking}. Se toma una estructura cristalográfica del receptor con su ligando, se extrae el ligando, se aleatoriza su conformación y se vuelve a acoplar con el protocolo que se quiere evaluar. Después se compara la pose obtenida con la cristalográfica mediante el RMSD de los átomos pesados, \emph{sin superponer} (ambas están ya en el sistema de referencia del receptor) y teniendo en cuenta la simetría del ligando (un anillo de benceno girado \ang{180} es la misma pose):
\begin{equation}\label{eq:15-rmsdsim}
\mathrm{RMSD}_{\mathrm{sim}}=\min_{\pi\in\mathrm{Aut}(G)}\sqrt{\frac1{N}\sum_{k=1}^{N}\big\lVert\mathbf p_k-\mathbf q_{\pi(k)}\big\rVert^2}.
\end{equation}

\begin{simbolos}
\simbolo{$\mathbf p_k,\ \mathbf q_k$}{Posiciones del átomo pesado $k$ del ligando en la pose predicha y en la cristalográfica.}
\simbolo{$\mathrm{Aut}(G)$}{Automorfismos del grafo molecular: permutaciones de átomos que dejan la molécula invariante.}
\simbolo{$N$}{Número de átomos pesados del ligando.}
\end{simbolos}

El criterio convencional es que una pose es correcta si $\mathrm{RMSD}_{\mathrm{sim}}<\SI{2}{\angstrom}$. Un protocolo que no reproduce la pose cristalográfica del propio ligando en su propio receptor difícilmente acertará con ligandos nuevos. Se informan dos tasas: la de éxito \emph{top-1} (la pose mejor puntuada es correcta) y la de éxito \emph{top-N} (alguna de las $N$ mejores lo es); la diferencia entre ambas separa los fallos de puntuación de los de búsqueda. Si la pose correcta nunca aparece entre las generadas, falló la búsqueda; si aparece pero no en primer lugar, falló la puntuación.

\begin{codigo}[redocking\_rmsd.py]
from rdkit import Chem
from rdkit.Chem import rdMolAlign

ref = Chem.MolFromMolFile("ligando_cristal.sdf")
poses = Chem.SDMolSupplier("poses_vina.sdf")
for i, pose in enumerate(poses, start=1):
    # CalcRMS no superpone y considera la simetría del ligando
    r = rdMolAlign.CalcRMS(pose, ref)
    ok = "correcta" if r < 2.0 else "incorrecta"
    print(f"pose {i}: RMSD = {r:5.2f} A  ({ok})")
\end{codigo}

\subsection{Limitaciones, y el aprendizaje profundo}

El \ingles{docking} clásico es útil pero tiene limitaciones que todo usuario debe conocer:

\begin{itemize}
  \item \textbf{Receptor rígido.} Las proteínas se ajustan al ligando (\ingles{induced fit}); una estructura cristalizada con otro ligando, o sin ninguno, puede tener el bolsillo cerrado o con cadenas laterales en posiciones incompatibles. Acoplar contra varias conformaciones del receptor (\ingles{ensemble docking}) o permitir flexibilidad en algunas cadenas laterales mitiga el problema a un coste computacional mayor.
  \item \textbf{Agua y protonación.} Las moléculas de agua que median puentes entre ligando y proteína, y los estados de protonación, rara vez se modelan bien.
  \item \textbf{Puntuaciones que no son afinidades.} Como vimos en el \cref{ej:vina}, el error en afinidad es del orden de 2 kcal/mol. En cribado virtual, las puntuaciones sirven para \emph{enriquecer} la lista de candidatos, no para predecir cuáles son activos.
  \item \textbf{Modelos predichos.} Acoplar contra un modelo de AlphaFold es posible, pero las cadenas laterales del bolsillo, incluso con pLDDT alto, pueden no estar en la conformación de unión; la exactitud suele caer respecto a las estructuras experimentales cocristalizadas.
\end{itemize}

El aprendizaje profundo también ha llegado al \ingles{docking}. \textcite{c15-corso2023} replantearon el problema como uno de \emph{modelado generativo}: en lugar de buscar el mínimo de una función de puntuación, DiffDock aprende a generar poses mediante un proceso de difusión sobre las traslaciones, rotaciones y torsiones del ligando, y un modelo de confianza ordena las poses generadas. En su evaluación sobre complejos de PDBBind, obtuvo una tasa de éxito \ingles{top-1} (RMSD $<\SI{2}{\angstrom}$) del \SI{38}{\percent}, frente al \SI{23}{\percent} de los mejores métodos de búsqueda tradicionales, sin necesidad de indicar el bolsillo. AlphaFold3, como vimos, va más allá y predice de una vez la estructura de la proteína con su ligando \parencite{c15-abramson2024}. Estos métodos no han jubilado a Vina: son más lentos por compuesto, pueden producir geometrías químicamente imposibles y su rendimiento cae en proteínas muy diferentes de las de entrenamiento. En la práctica actual, ambos enfoques se combinan.

\begin{cuidado}[Cinco comprobaciones antes de creer un resultado de docking]
\begin{enumerate}
  \item ¿Se validó el protocolo por \ingles{redocking} en la misma diana, con RMSD $<\SI{2}{\angstrom}$?
  \item ¿Es razonable la química de la pose (puentes de hidrógeno con geometría correcta, grupos polares no enterrados sin pareja)?
  \item ¿Son robustas las mejores poses frente a cambios en la semilla aleatoria, la \texttt{exhaustiveness} y el tamaño de la caja?
  \item ¿Se comparó la puntuación con la de ligandos conocidos y con la de señuelos (\ingles{decoys}) de propiedades similares?
  \item ¿Hay evidencia experimental independiente (mutagénesis, relación estructura-actividad) que apoye la pose?
\end{enumerate}
\end{cuidado}

\begin{ideaclave}
El \ingles{docking} es búsqueda más puntuación. Valide ambas con \ingles{redocking} ($\mathrm{RMSD}<\SI{2}{\angstrom}$) y use las puntuaciones para ordenar, no para medir afinidades.
\end{ideaclave}

\begin{resumen}
\begin{itemize}
  \item La conformación del esqueleto proteico se describe con dos diedros por residuo, $\phi$ y $\psi$, porque el enlace peptídico es plano. El diagrama de Ramachandran, predicho por estereoquímica en 1963, se confirma con estructuras reales: hélices en torno a $(-\ang{64},-\ang{40})$, láminas en torno a $(-\ang{110},\ang{120})$ y la glicina como excepción.
  \item La cristalografía, la RMN y la crio-EM miden señales distintas y producen modelos con medidas de calidad distintas (resolución, $R_{\text{free}}$, factor $B$, ensambles); en 2025 la crio-EM ya aporta el \SI{41}{\percent} de las nuevas entradas del PDB, cuyo formato oficial es mmCIF.
  \item El algoritmo de Kabsch obtiene la superposición óptima con una SVD de la matriz de covarianza cruzada, corrigiendo las reflexiones. El RMSD está dominado por los peores residuos y depende de la longitud; el TM-score mide la fracción de estructura que coincide y es comparable entre proteínas.
  \item La hipótesis de Anfinsen justifica predecir la estructura a partir de la secuencia; la paradoja de Levinthal se resuelve con paisajes de energía en embudo. La predicción pasó del modelado por homología a la coevolución (DCA) y al aprendizaje profundo.
  \item AlphaFold2 aprende de extremo a extremo con el Evoformer (atención sobre MSA y pares, actualizaciones triangulares) y un módulo de estructura equivariante, y estima su error con el pLDDT (local) y el PAE (relativo entre residuos). ESMFold sustituye el MSA por un modelo de lenguaje; AlphaFold3 extiende la predicción a complejos con ligandos y ácidos nucleicos.
  \item El \ingles{docking} combina una búsqueda en $6+N_{\mathrm{rot}}$ dimensiones con una función de puntuación aproximada, como la de Vina. Se valida por \ingles{redocking} con un RMSD simétrico menor de \SI{2}{\angstrom}, y sus puntuaciones sirven para ordenar poses, no para medir afinidades con precisión.
\end{itemize}
\end{resumen}

\begin{ejercicios}
\ej{1} Descargue \texttt{1UBQ} en formato mmCIF, calcule $(\phi,\psi)$ de todos sus residuos con el programa \archivo{phi\_psi.py} y dibuje el diagrama de Ramachandran. ¿Qué residuos tienen $\phi>0$? Compruebe cuántos de ellos son glicinas.
\ej{1} Usando la \cref{eq:15-bfactor}, calcule el desplazamiento tridimensional típico de un átomo con $B=\SI{30}{\angstrom^2}$ y con $B=\SI{120}{\angstrom^2}$. ¿Confiaría en la orientación de una cadena lateral de lisina cuyos átomos tienen $B\approx120$?
\ej{2} Implemente la función \py{kabsch_rmsd} y compruebe con estructuras sintéticas que: (a) si $\mathbf Q=\mathbf R_0\mathbf P+\mathbf t_0$ para una rotación conocida, recupera $\mathbf R_0$ y RMSD $=0$; (b) si $\mathbf Q$ es la imagen especular de $\mathbf P$, el RMSD no es cero cuando se aplica la corrección $d$, y sí lo es si se omite. ¿Qué significa esto biológicamente?
\ej{2} Demuestre que, para matrices ortogonales $\mathbf M$, se cumple $|M_{jj}|\le1$ y, por lo tanto, $\operatorname{tr}(\mathbf M\boldsymbol\Sigma)\le\sum_j\sigma_j$. ¿Por qué, cuando $\det(\mathbf V\mathbf U^\top)=-1$, la mejor rotación propia sacrifica precisamente el menor valor singular?
\ej{2} Calcule el TM-score entre la ubiquitina cristalográfica y cada modelo del ensamble de RMN \texttt{1D3Z}, con y sin la cola C-terminal. Compare la variación relativa del TM-score con la del RMSD y explique la diferencia a partir de la \cref{eq:15-tmscore}.
\ej{2} Obtenga de AlphaFold DB el modelo de una proteína de su interés que tenga estructura experimental publicada después de 2021. Superponga ambas en los residuos con pLDDT $>70$ y calcule RMSD y TM-score. ¿Coinciden las regiones de pLDDT bajo con las que faltan o son flexibles en la estructura experimental?
\ej{3} Implemente la función de puntuación de Vina (\cref{eq:15-dsup,eq:15-vina,eq:15-vinarot}) para un complejo proteína-ligando de su elección, usando solo los átomos pesados y una tipificación simplificada (C hidrofóbico, N/O donador o aceptor). Compare su valor para la pose cristalográfica con el que devuelve \py{Vina.score()}.
\ej{3} Elija cinco complejos proteína-ligando del PDB y realice un experimento de \ingles{redocking} con Vina variando \texttt{exhaustiveness} (1, 8, 32). Calcule las tasas de éxito \ingles{top-1} y \ingles{top-5} con RMSD simétrico $<\SI{2}{\angstrom}$, y decida en cada fallo si es un error de búsqueda o de puntuación.
\end{ejercicios}

\referenciascapitulo
